% Generated by scripts/assemble_book.py; edit shared papers instead.

\part{Constructions and statements}

\chapter{Cellular dessins and their infinite orders}
We start with the surface-order constructions of Schreiber \cite{SurfaceOrders,SurfaceDuality}, retaining the original sources unchanged. The generic-resolution source \cite{Persistence} is background, not a substitute for exactness of a particular complex. All primary algebra examples below come from explicit clean Belyi maps. Finite-dimensional objects occur as modules or specified quotient strata of those infinite algebras, not as replacements for them. The initial arithmetic target is a primitive nonprincipal Dirichlet $L$-function. Finite-order Hecke characters, Dedekind zeta functions and general Artin representations require separate comparison theorems; Artin holomorphy is not assumed.

A clean dessin on a compact oriented surface is specified by a nonempty finite dart set $H$, a fixed-point-free involution $\alpha$, and a permutation $\sigma$ such that $\langle\alpha,\sigma\rangle$ is transitive. Contract the white vertices of degree two. The resulting cellular graph has vertices the $\sigma$-orbits, edges the $\alpha$-orbits, and faces the $\phi$-orbits, where
\[
 \tau=\sigma^{-1},\qquad \phi=\alpha\sigma^{-1},\qquad
 2-2g=\#\operatorname{cyc}(\sigma)-|H|/2+\#\operatorname{cyc}(\phi).
\]
Write $[h]=\{h,\alpha h\}$. Our quiver has vertices these edges and arrows $a_h:[h]\to[\tau h]$. Paths are composed in travel order: $ab$ traverses $a$ first. Set
\[
 A=kQ/(a_h a_{\phi h}:h\in H),\quad J=A_{>0},\quad
 \Lambda=\varprojlim A/J^m.
\]
Each vertex has two incoming and two outgoing arrows. Each arrow has precisely one permitted successor $a_{\tau h}$ and one forbidden successor $a_{\phi h}$; these are distinct because $\alpha$ has no fixed points. The same holds for predecessors. This is the \emph{locally gentle} condition. Every permitted cycle repeats indefinitely, so $A$ is infinite dimensional. We reserve ``gentle'' without qualification for the convention additionally imposing finite dimension.

\begin{proposition}[The actual split order]\label{thm:order}
For a permitted cycle $c$ of length $\ell$, let $B_c$ be the hereditary cycle order in $M_\ell(k[[u_c]])$: diagonal and forward matrix entries are in $k[[u_c]]$, and entries crossing the chosen cycle cut are in $u_c k[[u_c]]$. Then $\Lambda$ is the suborder of $B=\prod_c B_c$ whose two diagonal constant terms labelled $h,\alpha h$ agree. With $R=k[[u]]$ acting by $u\mapsto(u_c I)_c$, it is finite free of rank $\sum_c\ell(c)^2$; the conductor contains $uB$.
\end{proposition}
\begin{proof}
In a cycle order a path of length less than $\ell$ is the corresponding matrix unit, multiplied by $u_c$ exactly when it crosses the cut; another full turn multiplies it by $u_c$. Thus completed paths in each permitted branch give all entries of $B_c$. An arrow belongs to only that branch, whereas $e_{[h]}$ has diagonal entries at both occurrences $h,\alpha h$. Hence the only condition on branchwise formal coefficients is the displayed equality of diagonal constants. This proves both injectivity and surjectivity of the path-normal-form identification. Multiplication respects it: switching branches across a paired edge is precisely a forbidden pair. The quotient $B/\Lambda$ is killed by $u$, so the submodule has the same generic rank as $B$. It is torsion free and finite over the DVR $R$, hence free of that rank. The path and $u$-adic topologies agree: $uA\subset J$, and if $L$ is the longest permitted cycle, any path of length at least $L+1$ is $u$ times a shorter nonempty path. Thus $J^{L+1}\subset uA$. This also justifies the completion used in the identification.
\end{proof}
The path-degree variable $q$ below is not $u$: one full turn has path degree $\ell(c)$ but $u$-degree one. Formal, Laurent and operator-norm completions will likewise remain distinct.

Here are five explicit global inputs. Permutations list darts starting at zero; $E=|H|/2$.
\begin{center}\small
\begin{tabular}{llll}
\toprule
Surface and Belyi map $\beta$ & $\sigma$ & $\alpha$ & $(g,\operatorname{rank}_R\Lambda)$\\\midrule
$\mathbb P^1$, $4z(1-z)$ & $(0)(1)$ & $(01)$ & $(0,2)$\\
$\mathbb P^1$, $(z+1)^2/(4z)$ & $(01)$ & $(01)$ & $(0,4)$\\
$y^2=x^3-1$, $x^3$ & $(012)(345)$ & $(03)(14)(25)$ & $(1,18)$\\
$\mathbb P^1$, $4z^3/(1+z^3)^2$ & $(012)(345)$ & $(03)(15)(24)$ & $(0,18)$\\
$y^2=x^3-x$, $x^2$ & $(0123)$ & $(02)(13)$ & $(1,16)$\\\bottomrule
\end{tabular}
\end{center}
The first is the nodal interval: $A=k[x,y]/(xy)$ and $\Lambda=k[[x,y]]/(xy)$. The second has $A=k\langle x,y\rangle/(x^2,y^2)$, still infinite dimensional. The two theta dessins have the same two permitted three-cycles but different face cycles, respectively $(051324)$ and $(04)(13)(25)$. The elliptic square has one permitted and one forbidden four-cycle. For example $x$ has pole order two on the cubic $y^2=x^3-1$; $x^3$ has ramification $(3,3)$ above zero, $(2,2,2)$ above one since $x^3-1=y^2$, and $(6)$ above infinity. The rational theta map has passport $(3,3),(2,2,2),(2,2,2)$. These ramification computations, and the analogous direct zeros/poles in the other rows, verify cellular Belyi inputs rather than merely abstract quivers.


\chapter{Projective resolutions, dual dessins and determinants}
Use right modules $P_i=e_iA$, $S_i=P_i/P_iJ$. Internal shift $\gshift{M}{n}$ places its old generator in degree $n$. Let $\mathcal R_n(i)$ be the \emph{indexing words} of length $n$ from $i$ with every adjacent pair forbidden, including $e_i$ for $n=0$. Such a word is usually zero in $A$; it still labels a summand.
\begin{theorem}[Exact resolutions]\label{thm:resolution}
There is a minimal graded projective resolution, with $P_0(i)=e_iA$ and for $n\ge1$,
\[
 P_n(i)=\bigoplus_{p=a_1\cdots a_n\in\mathcal R_n(i)}
 \gshift{e_{t(a_n)}A}{n},\qquad
 d_n(x)_p=a_nx\ \text{in the summand }a_1\cdots a_{n-1}.
\]
The augmentation $P_0(i)\to S_i$ is the quotient. For the completed order the identical finite-rank projective terms and formal-series maps resolve the residue simple. Ext and Tor against residue simples count forbidden words with the prescribed endpoints. In the graded case Ext has matching internal shift $n$.
\end{theorem}
\begin{proof}
The maps are right linear and degree zero. Adjacent differentials multiply by a forbidden pair. The augmentation kernel is the direct sum of $aA$ over initial arrows. For any arrow $a$, multiplication $e_{t(a)}A\to e_{s(a)}A$ kills exactly the permitted paths beginning with its forbidden successor $f(a)$. All other basis paths survive to distinct basis paths. Thus its kernel is $f(a)A$. For a fixed prefix, different last arrows have disjoint surviving supports; there is no cancellation between summands. This proves exactness inductively. For formal series the same assertion holds coefficient by coefficient and identifies the entire closed kernel, not just its dense polynomial subspace. Terms are idempotent projectives and all entries are arrows. Consequently applying Hom or tensor with a residue simple kills every differential, proving minimality and the Ext/Tor count. In $A$ minimality is graded; it does not require that $J$ be its ordinary Jacobson radical.
\end{proof}
The two arms never terminate in this closed class. Their labelled patterns follow $\phi$-cycles. An ungraded positive syzygy tail can have a shorter period by permutation of summands: the nodal interval has labelled period two but unordered tail period one; the elliptic theta has labelled period six and unordered tail period three. A simple need not be isomorphic to a positive syzygy; for the node it is $R$-torsion, whereas the positive syzygies are lattices.

Let $C_{ij}(q)=\sum_{d\ge0}\dim_k(e_iA_de_j)q^d$ and $E_{ij}(q)=\sum_{n\ge0}(-1)^n\dim_k\operatorname{Ext}_A^n(S_i,S_j)q^n$. Degreewise Euler characteristic of the resolution gives $EC=I$: every coefficient is finite, because a summand in homological degree $n$ starts in internal degree $n$. Constant terms are $I$, so the inverse is two-sided. Already signed $E(q)$ must not be evaluated at $-q$ again.

\begin{proposition}[Dual dessin and cycle determinant]\label{thm:cartan}
The complementary-relation algebra is another closed dessin algebra, with $\alpha^\#=\alpha$, $\sigma^\#=\phi^{-1}$, permitted successor $\phi$, and forbidden successor $\tau$. It exchanges graph vertices and faces and preserves genus. Its Cartan matrix satisfies $E(q)=C_{A^\#}(-q)$, with the opposite convention if required for Yoneda composition. For independent commuting arrow weights $w_a$,
\[
 \det C_A(\mathbf w)=
 \frac{\prod_{f\in\operatorname{cyc}(\phi)}(1-(-1)^{|f|}w(f))}
 {\prod_{c\in\operatorname{cyc}(\sigma)}(1-w(c))}.
\]
\end{proposition}
\begin{proof}
$\alpha(\sigma^\#)^{-1}=\alpha\phi=\tau$, so the claimed relations and face data follow directly; the connectedness group is unchanged. To prove the determinant, put $S_{v,a}=w_a$ when $s(a)=v$, and $T_{a,v}=1$ when $t(a)=v$, zero otherwise. On arrow states let $U_{a,b}=w_b$ for permitted adjacency and $V_{a,b}=w_b$ for forbidden adjacency. Then $TS=U+V$ and the path sum is $C=I+S(I-U)^{-1}T$. Sylvester's identity gives $\det C=\det(I+V)/\det(I-U)$. Each transition is a weighted permutation, whose cycle determinant is respectively $1-(-1)^\ell w(c)$ and $1-w(c)$. Everything holds in formal power series, hence as rational functions. This supplies a direct proof of the Bessenrodt--Holm formula in the closed dessin class \cite{BessenrodtHolm}.
\end{proof}
The five univariate determinants are, in the order above,
\[
 \frac{1+q}{1-q},\quad\frac{1+q}{1-q},\quad
 \frac{1-q^6}{(1-q^3)^2},\quad
 \frac{(1-q^2)^3}{(1-q^3)^2},\quad 1.
\]
The last cancellation does not imply finite global dimension: its forbidden resolution is infinite. These rational functions have cycle poles/zeros, not arithmetic prime Euler factors.

For clarity, a complex is a sequence with $d^2=0$, its homotopy category identifies maps differing by $dh+hd$, and $D(A)$ inverts maps inducing cohomology isomorphisms. A perfect complex is represented by a bounded complex of finite projectives. Our infinite simple resolutions compute derived Hom and tensor but are generally not perfect. The two-term presentation $[P_1\to P_0]$ is perfect even if its kernel is nonzero. Quadratic linearity is a Koszul statement; it is not by itself an equivalence between all completed derived categories. Likewise a finite-dimensional gentle graded-curve theorem \cite{GeometricGentle} cannot simply be transferred to these closed infinite orders. Here the surface walks explicitly compute the resolution endpoints: $\tau$ turns round a graph vertex and $\phi$ follows a face. For longer string objects all maps can be computed from the occurrence matrices; a universal geodesic-intersection classification for this completed category remains a separate theorem.


\chapter{Determinant semi-invariants and Schur constructions}
Assume $k$ algebraically closed of characteristic zero for invariant-ring statements. A finite-dimensional representation of the \emph{infinite} $A$ has spaces $N_i$, arrow matrices $N(a)$ satisfying the monomial relations, and base-change group $G=\prod_i\mathrm{GL}(N_i)$. With column vectors, $N(a)\mapsto g_{t(a)}N(a)g_{s(a)}^{-1}$. Let $P_0=\bigoplus P_i^{a_i}$, $P_1=\bigoplus P_i^{b_i}$, and $f:P_1\to P_0$ be a finite path-matrix map. When $\sum a_i\dim N_i=\sum b_i\dim N_i$, define
\[
 c_f(N)=\det D_f(N),\qquad D_f(N)=\operatorname{Hom}_A(f,N).
\]
It is intrinsically a determinant-line section; chosen bases give this scalar.
\begin{proposition}[Weights and exact homological scope]\label{thm:semi}
$c_f(gN)=\prod_i\det(g_i)^{b_i-a_i}c_f(N)$. Restricted determinants of this form span all semi-invariants on a closed $G$-stable relation or nilpotence locus. Set $M=\operatorname{coker}f$. On the square locus,
\[
 c_f(N)\ne0\quad\Longleftrightarrow\quad\operatorname{Hom}(M,N)=0
 \quad\Longleftrightarrow\quad\operatorname{Hom}(M,N)=\operatorname{Ext}^1(M,N)=0.
\]
If $f$ is injective, $\operatorname{coker}D_f=\operatorname{Ext}^1(M,N)$ and higher Ext vanish. Without injectivity neither that cokernel equality nor higher vanishing follows.
\end{proposition}
\begin{proof}
Naturality gives $D_f(gN)=G_1D_f(N)G_0^{-1}$, where $G_j$ repeats the relevant $g_i$. Take determinants. The contragredient coordinate-ring action has the inverse character. Schofield--van den Bergh's spanning theorem allows oriented cycles \cite[Theorem2.3]{SchofieldVdB}. Exactness of invariants for reductive $G$ in characteristic zero lifts a fixed-weight function from any closed stable locus to the ambient representation space; restrict the spanning determinants. Injectivity of the path map need not survive that restriction. Finally $\ker D_f=\operatorname{Hom}(M,N)$. Writing $K=\operatorname{im}f$, the exact sequence with projective $P_0$ expresses $\operatorname{Ext}^1(M,N)$ as $\operatorname{Hom}(K,N)/\operatorname{im}\operatorname{Hom}(P_0,N)$; the epimorphism $P_1\to K$ embeds this quotient in $\operatorname{coker}D_f$. A square finite matrix with zero kernel has zero cokernel. If $f$ is injective it is the actual length-one resolution, giving the stronger assertion.
\end{proof}
For the node, $f=\times x$, $M=\Lambda/(x)$, $N=\Lambda/(x,y^2)$ give $\dim\operatorname{coker}D_f=2$ but $\dim\operatorname{Ext}^1(M,N)=1$. For the elliptic square name arrows so that $a,b:0\to1$, $c,d:1\to0$, permitted cycle $acbd$ and forbidden cycle $adbc$. With $f=a$ and $N(a)=1$, $N(b)=N(c)=N(d)=0$, its Hom matrix is $[1]$, yet the subsequent forbidden resolution gives $\operatorname{Ext}^{2,3}(M,N)=k$. Thus nonvanishing is not all-degree derived orthogonality.

\begin{proposition}[What completion does to invariants]\label{thm:nilpotent}
A continuous $d$-dimensional representation of $\Lambda$ kills $J^d$. On the reduced nilpotent representation variety of fixed dimension vector every invariant regular function is constant. Nonzero-weight semi-invariants need not be constant.
\end{proposition}
\begin{proof}
Continuity with discrete finite-dimensional target kills some $J^m$; its Loewy filtration has at most $d$ nonzero strict steps. Choose vertexwise complements to this filtration. An arrow sends a layer $p$ into layers $q\ge p+1$. Scaling layer $p$ by $z^p$ conjugates each component by $z^{q-p}$. At $z=0$ the representation tends to the zero-arrow representation. Every invariant therefore has the same value at all points; reducedness and the Nullstellensatz give the ring assertion. In the square example the scalar $N(a)$ varies on a nilpotent locus and has weight $g_1/g_0$, so semi-invariants survive. Formal entries in a projective map are evaluated on this locus by truncating modulo $J^d$, a legitimate finite polynomial operation.
\end{proof}
On the node, $u=x+y$ is regular and $A$ is free of rank two over $k[u]$. The polynomial presentation $\times(u-\lambda)$ for $\lambda\ne0$ is injective and its determinant on $N$ is $\det(u_N-\lambda I)$, the characteristic polynomial of $u_N$ up to the sign $(-1)^d$. But $u-\lambda$ is a unit in $\Lambda$, with its geometric-series inverse, so those modules do not extend continuously to the completion. An $R$-free lattice and its $u$-nilpotent quotients are different categories. On a lattice a separately specified endomorphism $F$ has a meaningful relative characteristic polynomial over $R$; the scalar base action $uI$ alone does not create Frobenius.

\section{Young tableaux and a cyclic representation stratum}
For a nonempty partition $\lambda\vdash m$, choose a numbering $T$ of its boxes. On $V^{\otimes m}$ let $Y_T=R_TC_T$, the unnormalized row sum followed by the signed column sum. The classical identity $Y_T^2=h_\lambda Y_T$, where $h_\lambda$ is the hook product, makes $Y_T/h_\lambda$ an idempotent in characteristic zero. Its image is $S_\lambda(V)$, one Schur copy, rather than the whole symmetric-group isotypic component \cite{Etingof,SchurComplexes}. Its character is $s_\lambda$, computed by semistandard tableaux (weak rows, strict columns), and
\[
 \dim S_\lambda(k^d)=\prod_{(i,j)\in\lambda}\frac{d+j-i}{h_{ij}}.
\]
Since $f^{\otimes m}$ commutes with the place permutations, it induces $S_\lambda(f)$ and respects composition and zero maps. Applying this at every vertex and arrow therefore gives a genuine representation of the monomial dessin algebra: $S_\lambda(N(b))S_\lambda(N(a))=S_\lambda(0)=0$ for a forbidden pair. The empty partition does not have this property: it sends zero arrow maps to the identity of $k$.

The Littlewood--Richardson rule gives natural intertwiners
\[
 S_\lambda(V)\otimes S_\mu(V)\cong
 \bigoplus_\nu S_\nu(V)^{\oplus c_{\lambda\mu}^{\nu}}.
\]
Here $c_{\lambda\mu}^{\nu}$ counts semistandard tableaux of skew shape $\nu/\lambda$ and content $\mu$ whose row-reading word is Yamanouchi. Naturality makes the same decompositions intertwine all the arrows. The tensor operation is vertexwise with synchronous arrow maps, not $\otimes_A$ and not a Hopf-algebra tensor structure. There is no automatic tensor unit preserving these zero relations. These functors are not additive or exact; they cannot simply be applied degreewise to a projective resolution. They also need not preserve the fixed $R$-action on lattices.

In the elliptic theta set one of the two permitted three-cycle branches to zero. The other has arbitrary maps $X_1:V_0\to V_2$, $X_2:V_2\to V_1$, $X_3:V_1\to V_0$, with all spaces of dimension $d$. This is a closed representation stratum of our infinite dessin algebra. Put $M=X_3X_2X_1$.
\begin{theorem}[Cyclic invariants and a local positive kernel]\label{thm:schur}
The invariant ring on this stratum is $k[c_1(M),\ldots,c_d(M)]$, where $c_j$ are characteristic coefficients. Its total-arrow-degree Hilbert series is $\prod_{j=1}^d(1-t^{3j})^{-1}$. Common multidegree $(m,m,m)$ has basis $\chi_\lambda(M)=\tr S_\lambda(M)$, $\lambda\vdash m$, $\ell(\lambda)\le d$; other multidegrees have no invariants. Products obey the Littlewood--Richardson rule. Over $\C$, if $\|M_i\|\le R_0$ and $0\le r$ and $rR_0^2<1$, the kernel
\[
 K_r(M_i,M_j)=\det(I-rM_i\otimes\overline{M_j})^{-1}
 =\sum_\lambda r^{|\lambda|}\chi_\lambda(M_i)\overline{\chi_\lambda(M_j)}
\]
is absolutely convergent and positive semidefinite.
\end{theorem}
\begin{proof}
Cycle-trace generation for quivers, including relation quotients, is the Le Bruyn--Procesi theorem \cite{LeBruynProcesi,CrawYamagishi}. Every cycle here is a power of the full three-cycle, up to cyclic rotation. In characteristic zero Newton identities give the generators $c_j$. They are algebraically independent: take $X_1=X_2=I$ and arbitrary $X_3$. Scalar changes at the three vertices force equal arrow multidegrees. Restriction to that same section identifies degree $(m,m,m)$ with degree-$m$ symmetric polynomials in eigenvalues, whose Schur basis and product rule give the assertion. For the kernel use the ordinary Cauchy identity, first as a formal identity in eigenvalues. On Schur modules the operator norm of $S_\lambda(M)$ is at most $R_0^{|\lambda|}$, so its trace is bounded by $\dim S_\lambda(\C^d)R_0^{|\lambda|}$. The dimension Cauchy identity bounds the absolute series by $(1-rR_0^2)^{-d^2}$. Each summand is a rank-one Gram kernel, proving positivity by absolute summation.
\end{proof}
The term for the empty partition is a constant scalar feature, not an allowed module-valued tensor unit. In dimension two,
\[
 (\tr M)^2=\tr\operatorname{Sym}^2M+\det M,\qquad
 \tr\operatorname{Sym}^2M=\tfrac12((\tr M)^2+\tr M^2).
\]
More generally $\det(I-zM)=\sum_j(-z)^j\tr\bigwedge^jM$ and its reciprocal is $\sum_m z^m\tr\operatorname{Sym}^mM$. These are actual determinant/Schur connections, not identifications with an arithmetic $L$-function. On continuous finite-length $\Lambda$-modules the monodromy is nilpotent and all positive-degree characters vanish. The displayed Cauchy positivity is local representation-theoretic positivity, not the global Weil form.


\chapter{Lattice almost-split sequences and genuine diamond quotients}
For the interval dessin let $X=(x)$, $Y=(y)$ in $\Lambda=k[[x,y]]/(xy)$, $R=k[[u]]$, $u=x+y$. A lattice is a finitely generated $\Lambda$-module free over $R$.
\begin{theorem}[The nodal lattice calculation]\label{thm:nodear}
Every lattice is $\Lambda^r\oplus X^a\oplus Y^b$. The almost-split sequences in this lattice category are
\[
 0\to Y\lhook\joinrel\longrightarrow\Lambda\xrightarrow{\,\times x\,}X\to0,
 \qquad0\to X\lhook\joinrel\longrightarrow\Lambda\xrightarrow{\,\times y\,}Y\to0.
\]
In particular $\tau X=Y$, $\tau Y=X$. Multiplication by $u$ between two copies of either row, followed by reduction modulo $u$, gives an actual nonzero exact three-by-three diagram.
\end{theorem}
\begin{proof}
Normalize $M$ \emph{inside its total quotient module}: $\overline M=BM$, $B=k[[x]]\times k[[y]]$. The conductor gives $u\overline M\subset M\subset\overline M$. Both branch projections of $W=M/u\overline M\subset k^p\oplus k^q$ are onto. Intersect $W$ with the two factors; its remaining quotient identifies their two residual quotients. Separate bases for these pure factors and the common quotient put $W$ into pure first lines, pure second lines and diagonal pairs. Lift branch bases over the DVRs. These are respectively the summands $X,Y,\Lambda$. Their endomorphism rings are local.

Exactness of the first sequence is $\operatorname{ann}x=Y$. It does not split: a map $X\to\Lambda$ lands in $\operatorname{ann}y=X$, and its composite with $\times x$ is a nonunit in $\operatorname{End}X=k[[x]]$. Every map $\Lambda\to X$ factors through this epimorphism, $\operatorname{Hom}(Y,X)=0$, and every nonunit endomorphism of $X$ is divisible by $x$ and factors through $X\hookrightarrow\Lambda\to X$. Classification handles all finite sums; a unit $X$ component would be a retraction. This is right almost splitness. The dual argument for nonsections out of $Y$ proves left almost splitness. Endomorphisms fixing either structure map have constant term one, hence are units, proving minimality. Swap $x,y$ for the other sequence. Finally all three terms are $R$-free. Multiplication by $u$ is injective and reduction preserves exactness, proving the diagram assertion.
\end{proof}
The lattice AR quiver has arrows $Y\to\Lambda\to X$ and $X\to\Lambda\to Y$; the mesh composites are zero. Its projective vertex $\Lambda$ is not the endpoint of another nonsplit AR sequence. AR meshes do not mean that all rows and columns of any drawn grid are short exact.

For a second, explicit family fix $n\ge3$ and the rational Belyi map
\[
 \beta_n(z)=\frac1{1-T_n(z)^2}
 =\frac1{(1-z^2)U_{n-1}(z)^2},
\]
where $T_n,U_{n-1}$ are Chebyshev polynomials. It has passport $(2n),(2^n),(2^{n-1},1,1)$. Its contracted dual graph is a path, identifying its genus-zero dessin. The associated infinite algebra has vertices $1,\ldots,n$, endpoint loops $l,r$, opposite arrows $a_i:i\to i+1$, $b_i:i+1\to i$, and precisely the relations $l^2,r^2,a_ib_i,b_ia_i$. Its single permitted cycle is $l a_1\cdots a_{n-1}r b_{n-1}\cdots b_1$; the order rank is $(2n)^2$.

Retain one arrow at each edge in an orientation $\Omega$, and kill its opposite. The quotient is exactly the string algebra $H_n(\Omega)$ of Huang--Lin--Su \cite{HuangLinSu}; the killed-arrow ideal contains $J^{n+2}$, so it is also an actual quotient of $\Lambda_n$. Inflation identifies its module category with the full abelian quotient-annihilator subcategory $\mathcal C_\Omega$ of $\Lambda_n$-modules. This is generally not extension closed in the whole order category. The imported AR property is intrinsic to $\mathcal C_\Omega$, not automatically an ambient lattice or finite-length AR property. The extended edition gives nine explicit string modules, all twelve maps and the four exact meshes from the supplied diamond resource, now situated in this actual dessin quotient.

The other supplied resources also have a precise home. For the interval order the subcategory annihilated by $J^2$ is $\operatorname{Mod}B_0$, $B_0=\Lambda/J^2$. This is a domestic string quotient, not itself a gentle algebra. Inflation is full and preserves purity, and $\operatorname{Zg}B_0$ is a closed subspace of $\operatorname{Zg}\Lambda$. Ringel's infinite-word/pure-injective program and its proof by Puninski--Prest apply there \cite{Ringel,PuninskiPrest}; the isolated band $xy^{-1}$ gives Krull--Gabriel and relative Cantor--Bendixson dimension two by \cite[Theorem10.19]{LakingPrestPuninski}. This is not an ambient rank calculation for all $\Lambda$-modules. The parameter in $x=J_2,y=\lambda J_2$ survives even though all cycle traces vanish; the pp condition $\exists v\ (z=xv\ \wedge\ (y-\lambda x)v=0)$ detects it. These modules are $u$-torsion and are not nonzero order lattices. Infinite-string and locally finite AR theorems must retain their distinct finiteness hypotheses; a surviving permitted cycle has infinitely many paths at a vertex.

\section{Polynomial strings, completed lattices and the Fukaya category}
Crawley--Boevey's theorem \cite[Theorems1.1--1.4]{CrawleyBoeveyStrings} applies directly to the infinite polynomial dessin algebra: its finitely controlled modules decompose into strings and finite-dimensional bands; its pointwise artinian modules additionally allow primitive injective bands. For our finite quiver, these conditions mean finitely generated and artinian respectively. The words may be infinite, but their modules have algebraic direct-sum bases. This is not a classification of every module or of arbitrary formal completions.

In fact no nonzero completed-order lattice, restricted to $A$, meets either finiteness condition. The polynomial $A$ is finite over $k[u]$: remove full turns from sufficiently long paths as in Proposition~\ref{thm:order}. If an $R$-free lattice were finitely generated over $A$, it would be finite torsion free, hence free, over the PID $k[u]$. Multiplication by $1-u$ would then fail to be onto, whereas it is a unit over $R=k[[u]]$. Also $M\supsetneq uM\supsetneq u^2M\supsetneq\cdots$ excludes artinianity. For a concrete indecomposable outside the ordinary string/band list take $k[[x]]$ with $y=0$: its endomorphisms are multiplication by formal series, so its endomorphism ring is local. Its $x$ is injective and nonsurjective, while $1-x$ is invertible; these properties exclude every pure-$x$ ordinary string or band. The extended edition gives the full word argument.

The Fukaya connection requires equally precise categories. Lekili--Polishchuk explicitly allow infinite-dimensional locally gentle algebras, but their equivalence is for \emph{perfect dg modules} of a homologically smooth algebra, with specified grading, line field and stops \cite[Definition3.1.1, Lemma3.1.3, Theorem3.2.2]{LekiliPolishchuk}. Every nonempty closed dessin here has both permitted and forbidden cycles, so its polynomial algebra is neither proper nor smooth by their criterion. This does not exclude appropriate singular or completed variants; it prevents a direct use of that particular theorem. Infinite word, graded curve and geodesic are not interchangeable without a surface realization and boundary/grading data.

There is a genuine positive example inside our setting: the polynomial interval quotient $A/(y)=k[x]$ is still infinite dimensional, is smooth, and has no forbidden cycle. Its perfect category has the partially wrapped annulus model with one stop total. In degree zero $\operatorname{Perf}(k[x])=D^b(\operatorname{mod}_{\rm fg}k[x])$. Inflation of this quotient is fully faithful on ordinary modules, but not on derived categories: $\operatorname{Ext}_A^2(k[x],k[x])=k$, whereas $\operatorname{Ext}_{k[x]}^2(k[x],k[x])=0$. The former follows from the alternating $y,x$ resolution. Thus this quotient bridge does not identify the whole closed dessin's perfect category with the annulus category.

Finally, Leavitt path algebras \emph{can} fail invariant basis number, but need not. The permitted-state Laurent algebras here have finite-dimensional matrix quotients, and therefore invariant basis number. The polynomial and completed dessin algebras have their vertex-simple quotients, giving the same conclusion. The complex Leavitt algebra of $G_\alpha$ below is a dense *-subalgebra of its graph C*-algebra, on that specified graph; it is not the original gentle algebra with ordinary CK1 imposed.



\chapter{Leavitt relations, the constrained Fock algebra and spectrum}
Work now over $\C$. Ordinary Leavitt and graph C* constructions \cite{AbramsArandaPino,BPRS} cannot simply impose our original forbidden products: if $ab=0$ and CK1 says $a^*a=p_{t(a)}$, then $0=a^*ab=b$. Every arrow is a forbidden successor, hence every arrow and then every vertex is zero. CK2 is not needed. The ordinary graph-Toeplitz CK1 setting has the same collapse.

The valid permitted-arrow-state graph instead consists of the disjoint directed $\sigma$-cycles. Its Leavitt algebra is $\bigoplus_cM_{|c|}(k[z,z^{-1}])$, and its graph C* algebra is $\bigoplus_cM_{|c|}(C(\mathbb T))$. These forget the $\alpha$-gluing; both theta dessins give the same answer. They are not the completed original order. Nor is the ordinary Toeplitz algebra of a cycle of length $\ell>1$ equal to $M_\ell(\mathcal T)$: their $K_0$ ranks are $\ell$ and one respectively \cite{BrownloweLacaRobertsonSims}.

Let $\mathcal H=\ell^2(\mathcal P)$ for all permitted finite paths, including vacua $e_v$. Define $L_a\delta_p=\delta_{ap}$ when the prefix is permitted and composable, and zero otherwise. Let $P_v$ project onto paths starting at $v$, $Q_a=L_aL_a^*$ and $v_v$ onto the vacuum. Then
\[
 L_aL_{\phi(a)}=0,\quad
 L_a^*L_a=P_{t(a)}-Q_{\phi(a)}=v_{t(a)}+Q_{\tau(a)},\quad
 P_v-\sum_{s(a)=v}Q_a=v_v.
\]
These follow on basis vectors. The different initial projection avoids CK1 collapse. Write $\mathcal T_A=C^*(P_v,L_a)$, a constrained-path Fock C* algebra in the monomial/subproduct setting \cite{KakariadisShalit,ShalitSolel}. The polynomial $A$ acts faithfully: applying a finite path sum to the sum of vacua recovers its coefficients. This representation need not extend to all formal $\Lambda$ series: $\sum_{m\ge0}c^m$ along a cycle would send its vacuum to a nonsquare-summable sum.

\begin{theorem}[Compact extension and incidence index]\label{thm:fock}
Let $\mathcal H_v$ be the terminal-path sector. There is an exact sequence
\[
 0\longrightarrow\bigoplus_{v\in Q_0}\mathcal K(\mathcal H_v)
 \longrightarrow\mathcal T_A\longrightarrow
 \bigoplus_{c\in\operatorname{cyc}(\sigma)}M_{|c|}(C(\mathbb T))\longrightarrow0.
\]
With $\delta[u]=[1-W^*W]-[1-WW^*]$, its boundary map is
$\delta(c)=-\sum_{h\in c}e_{[h]}$, minus the unsigned incidence map $B_\Gamma$ of the cellular graph. Thus, for its $E$ edges,
\[
\begin{array}{c|cc}
 &K_0(\mathcal T_A)&K_1(\mathcal T_A)\\\hline
\Gamma\text{ bipartite}&\Z^{E+1}&\Z\\
\Gamma\text{ nonbipartite}&\Z^E\oplus\Z/2&0.
\end{array}
\]
\end{theorem}
\begin{proof}
All prefixes preserve the terminal sectors. For paths $p,q$ ending at $v$, $L_pv_vL_q^*=|\delta_p\rangle\langle\delta_q|$. Therefore the vacuum ideal is exactly the displayed sum of compacts: it contains these matrix units and that sum is itself an ideal. In the quotient, $L_a$ has initial projection $Q_{\tau(a)}$ and range $Q_a$, a graph CK family for the permitted-state cycles. The $Q_a$ have infinite rank and nonzero quotient images; the length gauge action descends. Gauge-invariant uniqueness \cite[Theorem2.1]{BPRS} identifies the quotient as claimed.

For a cycle $c$ based at arrow $a$, lift its corner unitary by $W=L_cQ_a+(1-Q_a)$. Then $W^*W=1$ and $1-WW^*=Q_a-L_cQ_aL_c^*$ projects onto the first $|c|$ paths starting with $a$. Their terminal vertices count exactly its incident darts, including twice for a loop. This proves the index formula. The six-term exact sequence gives $K_1=\ker B_\Gamma$ and $0\to\operatorname{coker}B_\Gamma\to K_0\to\Z^{\#\operatorname{cyc}\sigma}\to0$. The free last term splits abstractly. Spanning-tree integer elimination gives Smith entries $1$ repeated $V-1$ times and one zero in the bipartite case; an odd cycle instead adds the entry two. This proves the groups.
\end{proof}
The nodal interval has $K_0=\Z^2$, $K_1=\Z$; the single-loop sphere has $\Z\oplus\Z/2$, zero. Both theta dessins have $\Z^4$, $\Z$. Coarse $K$-theory does not recover genus or the arithmetic factors.

\section{A graph presentation and an actual pullback}
The Fock algebra itself has a graph presentation, on a different graph from the original gentle quiver. Let $G_\alpha$ have vertices $q_h$ for $h\in H$ and sinks $b_e$ for $e\in Q_0$, and edges
\[
 s_h:q_h\to q_{\tau h},\qquad t_h:q_h\to b_{[\tau h]}.
\]
\begin{proposition}\label{thm:graphmodel}
$C^*(G_\alpha)\cong\mathcal T_A$, by $q_h\mapsto Q_{a_h}$, $b_e\mapsto v_e$, $s_h\mapsto L_{a_h}Q_{a_{\tau h}}$ and $t_h\mapsto L_{a_h}v_{[\tau h]}$. The integer graph relation at $h$ is $q_h-q_{\tau h}-b_{[\tau h]}$. Eliminating these relations recovers $K_1=\ker B_\Gamma$ and $K_0\cong\operatorname{coker}B_\Gamma\oplus\Z^{\#\operatorname{cyc}\sigma}$.
\end{proposition}
\begin{proof}
The two edges at $q_h$ have orthogonal initial pieces and orthogonal ranges summing to $Q_{a_h}$, by the Fock initial-projection identity. All vertex projections are nonzero. The length gauge acts with degree one on both edge families, so gauge-invariant uniqueness applies. They generate since $L_{a_h}=s_h+t_h$ and $P_e=v_e+\sum_{[h]=e}Q_{a_h}$. In the integer matrix a kernel coefficient is constant around each permitted cycle; its sink equations are exactly $B_\Gamma c=0$. In the cokernel, eliminate all but one state per cycle; its remaining relation is the sum of its incident sink generators. The remaining cycle generators are free, proving the stated groups.
\end{proof}
For the nodal interval, $G_\alpha$ has two loop vertices, each with an exit to the common vacuum sink. Chirvasitu--Hajac--Tobolski's actual pullback theorem \cite[Theorem3.4]{GraphPullbacks} gives
\[
 \{T\in\mathcal T_A:\pi(T)\in\C^2\}
 \cong C^*(E_1),\qquad \C^2\hookrightarrow C(\mathbb T)^2\text{ by constants},
\]
where $E_1$ has the same three vertices and countably many arrows from each loop vertex to the sink, with no loops. The arrow numbered $n\ge0$ maps to loop$^n$ followed by exit. This algebra is AF, with $K_0=\Z^3$, $K_1=0$. The extended edition verifies every hypothesis. The theorem does not turn arbitrary cycle-vertex identifications into pullbacks: its first graph is acyclic and its path functor must cover all pointed paths. Its sink-extension Theorem4.3 has further injection and path-coverage requirements.

\section{A separate diagonal-fiber C* completion retaining the pairing}
The formal order's residue gluing cannot be copied by taking diagonals of arbitrary full matrices: $\operatorname{diag}(E_{12}E_{21})\ne\operatorname{diag}(E_{12})\operatorname{diag}(E_{21})$. For $\ell>1$ there is no nonzero *-homomorphism from $M_\ell(C(\mathbb T))$ to a commutative C*-algebra, as restriction to its constant matrix units shows. Restrict the distinguished fiber first. Define
\[
 D_\ell=\{f\in C(\mathbb T,M_\ell(\C)):f(1)\text{ is diagonal}\},\qquad
 \mathcal D_H=\Bigl(\prod_cD_{|c|}\Bigr)\mathop{\times}_{\C^H}\C^{Q_0},
\]
where the first map evaluates diagonal entries at one and the second repeats $\lambda_e$ at the two darts paired by $\alpha$.
\begin{theorem}[A faithful polynomial pullback model]\label{thm:diagonalmodel}
There is a split exact sequence
\[
 0\to\bigoplus_c C_0(\mathbb T\setminus\{1\},M_{|c|})
 \to\mathcal D_H\to\C^{Q_0}\to0,
\]
so $K_0(\mathcal D_H)=\Z^{|Q_0|}$ and $K_1(\mathcal D_H)=\Z^{\#\operatorname{cyc}\sigma}$. The polynomial dessin algebra embeds faithfully and generates this C*-algebra with its adjoints. This embedding does not extend to a unital algebra map from all of $\Lambda$.
\end{theorem}
\begin{proof}
The section consists of the constant paired diagonal projections. The kernel is the displayed vanishing fiber, a sum of matrix suspensions, giving the groups by split exactness and suspension. For an explicit embedding set $g(z)=|1-z|^2/4$, choose one cut arrow per permitted cycle, and send $a_h$ to $g(z)z^{\epsilon_h}E_{h,\tau h}$ in its branch, where $\epsilon_h=1$ at the cut and zero otherwise. Send $e_e$ to the constant paired diagonal projection. Every forbidden product is zero. A full cycle at $h$ maps to $w_c(z)E_{hh}$, $w_c=g^{|c|}z$. This scalar function is injective on the circle: away from its unique zero at one its argument recovers $z$. Continuous functional calculus therefore supplies every $f(z)E_{hh}$ with $f(1)=0$. Multiplying by path images and approximating division away from one supplies every vanishing off-diagonal entry. Together with the constant paired diagonals these generate $\mathcal D_H$. Path normal forms in a fixed entry are a nonzero prefactor times a polynomial in $w_c$; its infinite range proves independence. Evaluation at one first separates the zero-length coefficients. Finally for $z_0\ne1$, $1-w_c(z_0)^{-1}c_h$ is a unit in the formal order, but its proposed image is singular at $z_0$. A unital map preserves units, proving nonextension.
\end{proof}
This is the C* algebra generated by a faithful polynomial representation and its adjoints, distinct from the completed order and from $\mathcal T_A$. The norm closure of the polynomial image alone need not contain those adjoints. Their $K$ groups already differ. The pairing belongs to the labelled presentation; an unmarked isomorphism class need not recover the ribbon ordering. No arithmetic descent, Frobenius action or completed-$L$ positive determinant is inferred from this additional valid gluing model.


\begin{proposition}[The degree spectral triple]\label{thm:degree}
Define $N\delta_p=|p|\delta_p$ on $\{\xi:\sum_p|p|^2|\xi_p|^2<\infty\}$. With the polynomial creation/annihilation *-algebra this is an ungraded spectral triple, $p$-summable exactly for $p>1$. With $N_+^{-s}$ zero on the kernel,
\[
 \operatorname{Tr}e^{-tN}=E+\frac{2E}{e^t-1},\quad
 \operatorname{Tr}N_+^{-s}=2E\zeta(s),\quad
 \operatorname{Tr}\bigl(\chi(N)N_+^{-s}\bigr)=2E L(s,\chi)
\]
for $t>0$ and $\Rea s>1$ in the last two formulas, where $\chi$ is an externally chosen Dirichlet character.
\end{proposition}
\begin{proof}
There are $E$ vacua and $2E$ paths of every positive length: each starting arrow has a unique continuation. The real diagonal operator is self-adjoint on the stated maximal domain; finite supports are a core and its resolvent is compact. Prefixes preserve the domain and $[N,L_a]=L_a$, $[N,L_a^*]=-L_a^*$, so polynomial commutators are bounded. The trace formulas and summability reduce to the corresponding scalar series. The character multiplier is bounded; membership in $\mathcal T_A$ is not asserted.
\end{proof}
Length parity anticommutes with the creators but commutes with $N$, so it is not an even-spectral-triple grading. The ordinary odd Fredholm class of this positive $N$ is trivial: $N(1+N^2)^{-1/2}-I$ is compact. The spectrum depends only on $E$; zeros of the scalar zeta trace are not eigenvalues of $N$.


\chapter{Arithmetic attachment, descent and the remaining comparison}
The compact ideal gives a precise home for an arithmetic determinant, with no claim that the arithmetic is derived from the dessin. Fix one terminal sector. Its two paths of each length $n\ge1$ give orthonormal vectors $\xi_n=(\delta_{p_{n,1}}+\delta_{p_{n,2}})/\sqrt2$ and an isometry $U:\ell^2(\mathbb N_{\ge1})\to\mathcal H_v$.

For a nonprincipal character $\chi$ put $A_n=\sum_{j\le n}\chi(j)$, $B_0=\sup|A_n|\le q_\chi$, and, on $2^k\le n<2^{k+1}$, set
\[
 w_n=\sqrt n(k+1),\quad u_n=A_n/w_n,\quad
 v_n(s)=w_n(n^{-s}-(n+1)^{-s}),\quad
 \varphi_s(z)=\sum_n v_n(s)z_n-z_1.
\]
Here $u_1=1$. For $\Rea s>0$, $u,v(s)\in\ell^2$ and the pairing is complex linear, with no conjugation of $s$. Define
\[
 \widetilde K_s=U(u\otimes\varphi_s)U^*,\qquad
 \det_F(I+\widetilde K_s)=L(s,\chi).
\]
Indeed dyadic blocks give $\|u\|^2\le2B_0^2$, and on a compact set $\Rea s\ge\sigma_0>0$, $|s|\le M$, the squared $v$ tail is bounded by $M^2\sum_{k\ge J}(k+1)^2 2^{-2\sigma_0k}$. Thus the rank-one family is trace-norm holomorphic. Its determinant is $1+\varphi_s(u)$, equal by Abel summation to $L$. Replacing $v$ by $P_Nv$ and adding $A_N(N+1)^{-s}e_1^*$ to the functional gives exactly every raw partial sum $S_N$, with compact-uniform trace-norm convergence; $\|UXU^*\|_1=\|X\|_1$. For primitive $\chi$, completion replaces the functional by $g(s)\sum v_n(s)z_n-z_1$, with the standard nonzero gamma factor $g$, rather than multiplying the whole old functional. The extended edition supplies the quantitative estimates.

This inserts a scalar encoding into one actual dessin compact ideal; it is not known to be $A$-linear, a derived chain map or Frobenius compatible. A common fixed positive metric cannot make this same fixed-range family symmetric for all critical heights: symmetry forces its functionals to be proportional, whereas the coefficient ratio at heights $0,t\ne0$ is asymptotic to $(\sigma+it)n^{-it}/\sigma$, with two opposite subsequential limits. Different positive realizations remain possible.

\begin{proposition}[Complex dessin data do not determine arithmetic descent]\label{thm:descent}
The rational Belyi pairs $E_d:y^2=x^3-d$, $\beta_d=x^3/d$, for $d=1,2$, are isomorphic over $\C$ but have distinct Frobenius polynomials at the good prime seven: $1-4T+7T^2$ and $1-T+7T^2$.
\end{proposition}
\begin{proof}
Choose $a^3=d$, $b^2=d$. The projective isomorphism $(x,y)\mapsto(ax,by)$ from $E_1$ to $E_d$ preserves the Belyi map. The degree-six clean passports are the elliptic theta passport already computed. Directly counting $y$ for $x=0,\ldots,6$ gives respectively $(0,1,1,0,1,0,0)$ and $(0,0,0,2,0,2,2)$. Including infinity yields four and seven points, hence $a_7=8-\#E_d(\mathbb F_7)=4,1$. Both discriminants $-432d^2$ are nonzero modulo seven.
\end{proof}
Any invariant of the complex dessin or its undecorated order/Fock algebra alone is the same for these two inputs, so cannot recover both factors without descent data. This does not obstruct a theory including a specified rational model and compatible Galois action. Elliptic $H^1$ has weight one and unshifted center $s=1$, unlike Dirichlet center $1/2$; the example is not a counterexample to GRH.

\section{A constructive arithmetic benchmark on the theta dessin}
The preceding twists can be distinguished constructively once their finite-field descent data are retained. This closes the elementary local-factor computation, not the global positive-determinant problem. Let $\kappa=\mathbb F_q$, $\operatorname{char}\kappa\ne2,3$ and $d\in\kappa^*$. Put $E_d:y^2=x^3-d$ and $a_q(d)=q+1-\#E_d(\kappa)$. For $q\equiv1\pmod3$, fix a nontrivial cubic character $\rho$ and the quadratic character $\eta$, both extended by zero at zero. Define
\[
 J=\sum_{u\in\kappa}\rho(u)\eta(1-u),\qquad
 j=-\eta(-1)J,\qquad \lambda_d=\eta(d)\rho(d),\quad \pi_d=\lambda_dj.
\]
\begin{theorem}[Uniform local factor]\label{thm:jacobitheta}
For every such prime power and nonzero $d$,
\[
 a_q(d)=\pi_d+\overline{\pi_d},\quad |\pi_d|^2=q,\quad
 P_d(T):=1-a_q(d)T+qT^2=(1-\pi_dT)(1-\overline{\pi_d}T).
\]
If $q\equiv2\pmod3$, instead $a_q(d)=0$. The split-case values lie in $\Z[\omega]$, $\omega^2+\omega+1=0$.
\end{theorem}
\begin{proof}
Cube fibers have size $1+\rho(t)+\overline{\rho(t)}$, including at zero; the first term is the literal constant one, not the zero-extended trivial character. Counting square fibers gives $a_q(d)=-\sum_x\eta(x^3-d)$. Grouping by $t=x^3$ removes the constant term since $\sum_t\eta(t-d)=0$. Substitution $t=du$ gives $\sum_t\rho(t)\eta(t-d)=\rho(d)\eta(-d)J$; the other term is its conjugate. This proves the signed trace formula.

For a nontrivial additive character $\psi$, write $G(\chi)=\sum_x\chi(x)\psi(x)$. Character orthogonality gives $|G(\chi)|^2=q$ for nontrivial $\chi$: substitute $x=uy$ in its squared norm; the inner sum over $y\ne0$ is $q-1$ at $u=1$ and $-1$ otherwise. Grouping $G(\rho)G(\eta)$ by $x+y$ gives $JG(\rho\eta)$, the zero-sum fiber vanishing since $\rho\eta$ is nontrivial. The three Gauss norms yield $|J|^2=q$. Neither purity nor RH was assumed. The remaining signs and character values have modulus one. For $q\equiv2\pmod3$ cubing is bijective, so the quadratic-character sum is zero. Finally all split-case summands are signs times cube roots of unity. These are the classical Gauss--Jacobi methods \cite[pp.499--502]{WeilJacobi}, with our normalization proved explicitly.
\end{proof}
At $q=7$, choose $\rho(3)=\omega$. Then $j=3+2\omega$, $\lambda_2=\omega^2$, and $\pi_2=-1-3\omega$, giving the two factors already counted. At $q=13$, $\rho(2)=\omega$ gives $J=1+4\omega$ but $j=-1-4\omega$; this detects the essential sign hidden by the first example.

\section{A genuine cellular carrier, and what still fails for the order}
Assume $q\equiv1\pmod3$ in this cellular construction. The elliptic theta dessin has cellular chains $C_2=\mathbb{Q}$, $C_1=\mathbb{Q}^3$, $C_0=\mathbb{Q}^2$. Orient the three edges from $v_+$ to $v_-$; then $\partial_1(e_i)=v_--v_+$ and $\partial_2=0$. Thus $H_1$ is the sum-zero plane, with basis $e_0-e_1,e_1-e_2$. Rotation $\sigma$ acts there by
\[
 C=\begin{pmatrix}0&-1\\1&-1\end{pmatrix},\qquad
 G=\begin{pmatrix}2&-1\\-1&2\end{pmatrix}.
\]
$G$ is the actual Euclidean edge Gram matrix in this basis. The deck involution $\alpha$ acts as $-I$. Consequently the six-element deck group $\langle\sigma,\alpha\rangle$ acts on this homology and permutes the original quiver and its entire projective resolution.

Write $\pi_d=A_d+B_d\omega$. Its multiplication action on this \emph{actual cellular} plane is
\[
 M_d=A_dI+B_dC=\begin{pmatrix}A_d&-B_d\\B_d&A_d-B_d\end{pmatrix},\qquad
 \det(I-TM_d)=P_d(T),\qquad M_d^tGM_d=qG.
\]
This constructs the local matrix from finite-field character sums and the dessin's deck action, rather than fitting its trace. Positivity follows from $2x^2-2xy+2y^2=(x-y)^2+x^2+y^2$. The normalized matrix is $G$-unitary, not asserted self-adjoint. The extended edition supplies a rational cellular chain lift, its exact determinant cancellation, and Kummer descent compatibility on homology. Its odd-over-even determinant ratio is the classical finite-field curve zeta function
\[
 Z(E_d/\kappa,T)=\frac{P_d(T)}{(1-T)(1-qT)}.
\]
Identification with the geometric $\ell$-adic Frobenius polynomial uses the standard elliptic trace and determinant theorem; a canonical integral comparison or an order-derived Frobenius functor has not been constructed here.

There is also a precise obstruction to extending this same real two-dimensional carrier to $q\equiv2\pmod3$. Frobenius then conjugates the cubic deck action. The relation $MC=C^2M$ forces $M=\left(\begin{smallmatrix}a&b\\a+b&-a\end{smallmatrix}\right)$, whose determinant is $-a^2-ab-b^2\le0$, contradicting the required positive determinant $q$. This excludes that real carrier with its fixed deck action, not complex or enlarged Hilbert--P\'olya constructions. The extended edition gives a compatible complex local replacement.

For the following insertion test use the characteristic-zero theta order $\Lambda_{\mathbb{Q}}$ over $R_{\mathbb{Q}}=\mathbb{Q}[[u]]$, with residue simple $S_e=\mathbb{Q}$. The limitation is concrete. The deck action has finite order and cannot alone supply the norm-$\sqrt q$ Jacobi coefficient. Inserting $M_d$ into $S_e\otimes\mathbb{Q}^2$ gives an actual residue-module endomorphism with determinant $P_d$, but that simple is not perfect over the order and choosing it breaks full deck symmetry. Inserting the same matrix $\Lambda_{\mathbb{Q}}$-linearly into the regular projective instead gives the $R_{\mathbb{Q}}$-determinant $P_d^{18}$, or $P_d^6$ on $e\Lambda_{\mathbb{Q}}\otimes\mathbb{Q}^2$. Taking a root is not an established determinant-line comparison; a deck-semilinear lift is a different map. These surviving multiplicity/equivariance obligations prevent promoting the local computation to the required arithmetic order functor.

For the global comparison fix a nonzero rational $d$ and use its good reductions. The finite-field curve zeta has center $1/2$ after $T=q^{-s}$. The \emph{global elliptic} $L$-function over $\mathbb{Q}$ has unshifted center one, whereas Dirichlet $L$-functions have center $1/2$. Local weight-one purity does not determine either global zero set. The elementary local bound proves absolute convergence of the elliptic Euler product on $\Rea s>3/2$; the new local identities give no determinant-preserving continuation with the desired global positive spectrum.


\section{Arbitrary projective actions: an obstruction and a genuine determinant repair}
The preceding tensor test has a stronger replacement. Keep the characteristic-zero completed elliptic theta order $\Lambda$ over $R=\mathbb Q[[u]]$, and write $K=\mathbb Q((u))$. All determinants in the next theorem are \emph{ordinary underlying-$R$ determinants}; they are not reduced norms.
\begin{theorem}[Projective determinant multiplicity]\label{thm:theta-multiplicity}
Let $P$ be any finitely generated projective right $\Lambda$-module and $f\in\operatorname{End}_\Lambda(P)$. Then $P\simeq\bigoplus_{i=0}^2(e_i\Lambda)^{n_i}$. Let $A_i\in M_{n_i}(\mathbb Q)$ be the induced maps on its three top multiplicity spaces. Then
\[
 \left.\det_R(1-Tf)\right|_{u=0}
 =\left(\prod_{i=0}^2\det(1-TA_i)\right)^6.
\]
For any bounded complex of these projectives and any $\Lambda$-linear chain endomorphism, its alternating determinant ratio specializes to a sixth power in $\mathbb Q(T)$. Generically, the ordinary ratio is a cube in $K(T)$.
\end{theorem}
\begin{proof}
The arrow ideal is the Jacobson radical: its elements admit geometric-series inverses after subtracting from one, and the quotient is $\mathbb Q^3$. A projective's top lifts to a map from the displayed sum; Nakayama gives surjectivity, projectivity splits it, and another application of Nakayama kills the kernel.

Put $B=\Lambda/u\Lambda$. Each $e_iB$ has path layers of dimensions $1,2,2,1$: the identity, two length-one paths, two length-two paths, and one length-three path. The two full cycles based at $i$ add to $ue_i$, so become negatives; all length-four paths vanish. A positive-length entry of a projective homomorphism strictly raises path length. Consequently the associated graded map has six copies of $A=\operatorname{diag}(A_0,A_1,A_2)$, giving the formula by triangular block determinants. Base change of determinants justifies specialization. Multiplicativity along filtrations is ordinary linear algebra, consistent with \cite[Tags0FJ9,0GSX]{StacksDeterminants}; no purity assumption enters.

Each term's determinant has constant term one in $T$, so every denominator has nonzero specialization. Taking the finite alternating product therefore preserves the sixth power, even after cancellation. Finally $\Lambda_K=M_3(K)\times M_3(K)$; a projective's branch modules are sums of three-dimensional simples. A commuting endomorphism acts on each multiplicity space, repeating every factor three times. This proves the generic assertion.
\end{proof}
Thus no such ordinary perfect-complex ratio can equal the constant-in-$u$ quadratic $P_d(T)$, or the full local curve zeta $P_d(T)/((1-T)(1-qT))$. For $P_d$, the valuation at infinity is $-2$, not a multiple of six. For the zeta ratio, $P_d(1)=q+1-a_q>0$ by the local bound $|a_q|\le2\sqrt q<q+1$, so its pole at $T=1$ is simple. This also covers square prime powers where the numerator may have a double root. It is an obstruction for the specified ordinary determinant and linear category, not all spectral or determinant theories.

There is nevertheless a \emph{genuine integral functor} removing this multiplicity. Choose a permitted branch and its cycle cut. Its projection $\Lambda\to B_+\subset M_3(R)$ makes the column lattice $W_+=R^3$ a left $\Lambda$-module. Define
\[
 \Phi_+ :\operatorname{Perf}(\Lambda)\longrightarrow\operatorname{Perf}(R),
 \qquad P^\bullet\longmapsto P^\bullet\otimes_\Lambda W_+.
\]
Bounded projectives compute this derived tensor functor; flatness of $W_+$ on all modules is unnecessary. Since $e_iW_+\simeq R$, its determinant on $e_i\Lambda\otimes_{\mathbb Q}V$ with external action $1\otimes M_d$ is exactly $P_d$, whereas on the regular module it is $P_d^3$. This is a defined change of determinant functor, not an arbitrary polynomial root. It does not construct the arithmetic action independently: the already computed $M_d$ is still the input.

The price of this repair is precise. Branch swap $\alpha$ exchanges $W_+$ and $W_-$; pairing them yields $P_d^2$ on that basic-projective input. The chosen branch determinant cannot intertwine $\alpha$ on all endomorphisms: the corner unit $e_i+c_i^+$ acting on $e_i\Lambda$, with $c_i^+$ a plus full cycle, acts on its determinant by $1+u$, whereas its $\alpha$-image acts by one. Also no rank-one-per-vertex column lattice supported on one branch admits an $R$-linear lift of the rotation $\sigma$ fixing $u$. Its three nonzero arrow scalars have nonnegative $u$-valuations summing to one; an invertible cyclic lift would make all three equal. This contradiction concerns that lattice, not every determinant-level rotation comparison. Ramification $u=v^3$ \emph{and a new lattice} with all three arrow scalars $v$ admits such a rotation, changing the stipulated base/lattice data. A descent-compatible arithmetic determinant retaining the required symmetry and a global positive realization remains open.


\section{Tame descent: a surviving order and a sharper obstruction}
We now test the proposed ramified repair rather than assume that its rotation descends. Extend coefficients explicitly to $k=\mathbb Q(\omega)$, $\omega^2+\omega+1=0$. Put $R=k[[u]]$, $S=k[[v]]$, $u=v^3$, and let $g(v)=\omega v$, fixing $k$. On the balanced plus column $W=S^3$, let $C\epsilon_i=\epsilon_{i+1}$ and
\[
 \Theta=Dg,\qquad D=\operatorname{diag}(1,\omega^2,\omega).
\]
The plus arrows act by $vE_{i,i-1}$, and minus arrows by zero. Direct calculation gives $\Theta^3=1$, compatibility of $\Theta$ with this $\Lambda_R$-action, and
\[
 \Theta C=\omega^2C\Theta,\qquad
 L=W^\Theta=R\epsilon_0\oplus Rv\epsilon_1\oplus Rv^2\epsilon_2.
\]
Thus invariants restore precisely the original cut lattice. The extension $R\to S$ is finite \emph{free} of rank three and faithfully flat; the group quotient is not a torsor at $v=0$. The comparison $S\otimes_R L\to W$ has matrix $\operatorname{diag}(1,v,v^2)$, determinant $u$ and cokernel $S/(v)\oplus S/(v^2)$. Hence this semilinear group action alone is not effective descent of the column to $R$. This is the explicit local root-stack situation, consistent with the trivial-stabilizer criterion for vector-bundle descent \cite[Theorem10.3]{AlperDescent}; no stack theorem is needed for the displayed calculation.

Two positive constructions survive. The determinant line $\bigwedge_S^3W$ has trivial $D$ and $C$ actions and does descend, but the natural map
\[
 \det_R(L)\longrightarrow(\det_S W)^\Theta
\]
has image $uR$ inside $R$, not an isomorphism. Determinant and coarse invariants cannot be interchanged while dropping this nonunit comparison. Also the \emph{carrier} endomorphism algebra $\operatorname{End}_S(W)^\Theta$, conjugated to the invariant basis, is exactly the hereditary order with entries $R$ on and below the diagonal and $uR$ above. Its rotation descends as $\operatorname{Ad}(A)$, where
\[
 A=\begin{pmatrix}0&0&u\\1&0&0\\0&1&0\end{pmatrix},\qquad A^3=uI.
\]
The algebra automorphism has order three although $A$ is not an integral unit. This is not $\operatorname{End}_\Lambda(W)$ or a descended arithmetic cohomology group.

The natural arithmetic attempt fails: $aI+bC$ does not commute with $\Theta$ when $b\ne0$, as in both prime-seven Jacobi examples. The original cellular sum-zero plane is not $\Theta$-stable either. Replacing $C$ by $vC$ does give the descended operator $A$ on $L$, but
\[
 \det(1-T(aI+bA))=(1-aT)^3-ub^3T^3,
\]
whose special fiber is $(1-aT)^3$, not the required local quadratic. These carrier maps are not automatically $\Lambda$-linear arithmetic endomorphisms.

There is a stronger category-level conclusion.
\begin{theorem}[Full-deck coefficient kernels]\label{thm:full-deck-kernel}
Let $W_0$ be a finite $R$-free left $\Lambda_R$-module on which the central $u$ acts by the given scalar. Suppose it has invertible $R$-linear intertwiners lifting the order automorphisms $\sigma$ and $\alpha$. Then, for some $m\ge0$, every $e_iW_0$ has rank $6m$. For any finite projective $P$ and $\Lambda_R$-linear $f:P\to P$, let $A_i$ be its top blocks. The genuine coefficient functor $\Phi_{W_0}= -\otimes_\Lambda^{\mathbf L} W_0$ satisfies
\[
 \left.\det_R(1-T\Phi_{W_0}(f))\right|_{u=0}
 =\left(\prod_i\det_k(1-TA_i)\right)^{6m}.
\]
Bounded perfect ratios retain this exponent. Thus this whole class of full-deck, fixed-center finite free coefficient kernels cannot supply either the single local quadratic or the full local curve zeta.
\end{theorem}
\begin{proof}
In the generic fiber intersect $W_0$ with its two central branch summands to obtain finite free integral lattices $W_\pm$; an integral direct-sum decomposition is not assumed. Rotation preserves them and branch swap exchanges them. On a branch let the common vertex rank be $r$. The three arrow matrices have cyclic product $uI_r$. Their determinant valuations are nonnegative, sum to $r$, and rotation intertwiners make them equal. Therefore $3\mid r$. Branch swap makes the two ranks equal, so the full vertex rank is $6m$.

For the determinant formula, filter $W_0/uW_0$ by the arrow radical, whose fourth power vanishes. A positive-length entry of $f$ raises the filtration; a top block $A_i$ repeats across its $i$-th graded pieces a total of $\dim_k e_i(W_0/uW_0)=6m$ times. Triangular determinants give the assertion. Finite alternating products preserve it. Quadratic degree and the uncancelled simple pole of the full local zeta exclude the targets as before.
\end{proof}
This theorem is sharp: $W_0=\Lambda_R$ has vertex rank six. It requires neither left projectivity nor left flatness of a general kernel. It excludes this coefficient category, not arbitrary semilinear correspondences, twisted determinant lines or infinite spectral constructions. The extended edition identifies the full character-orbit repair with the hereditary normalization and explains why Schur commutator cancellation is insufficient. All determinant multiplicities here are relative to the stated $k$-valued base; restricting a rational quadratic factor to the former $\mathbb Q$-base squares it. No arithmetic descent or global positive determinant has been recovered by this ramification.


The recent NCG literature supplies concrete comparison targets, with strong hypotheses. Connes--Consani construct finite abelian covers of the adelic Riemann sector with prime-orbit length $\log p$ and arithmetic Frobenius monodromy \cite[Theorem3.12]{ConnesConsaniKnots}. The prolate/local-factor and moment constructions \cite{Prolate,LocalMoments} give particular finite-prime analytic models, not a uniform positive comparison for all completed $L$-functions. The zeta-spectral-triple determinant theorem \cite[Theorem5.10]{ZetaTriples} assumes a specified restricted Weil form, simple even bottom eigenfunction and boundary normalization. Its infinite comparison is still an additional step. The distributional real-zero theorem \cite[Theorem6.1]{QuadraticForms} concerns a particular interval quadratic form, essential self-adjointness on its trigonometric core and an even simple ground state; a Fock vacuum alone verifies none of those arithmetic identifications. Deninger's periodic-packet construction and real-cohomology obstruction \cite{DeningerPackets,DeningerReal} further specify, rather than remove, these issues. Full theorem-level reading coverage and exact versions are recorded alongside this draft.

A complete approach must construct arithmetic descent and prime actions on actual order/derived or analytic objects, prove determinant-line and local-factor compatibility, control regularization and conductor/gamma normalization, and establish a positive realization for that \emph{same} function. Finite-field cohomological trace formulas and Deligne purity supply the classical model \cite{DeligneWeilI,DeligneWeilII}; alternating duality alone is not positivity or purity. In a number-field construction these are obligations, not imported conclusions.

For a sharply testable final analytic target, put $F_\chi(z)=\Lambda(1/2+iz,\chi)$. Suppose entire functions $h_m(z)\det(zI-H_m)$, with $H_m$ finite self-adjoint and $h_m$ entire nowhere zero, converge locally uniformly on $\C$ to this nonzero completed function. Hurwitz on every disk disjoint from $\R$ then proves all its zeros real, hence the Dirichlet critical-line assertion. Establishing this convergence with independently constructed $H_m$ is exactly the missing load-bearing claim; a rank-one encoding or a reformulation does not establish it. Hecke, Dedekind and Artin extensions must state their analytic scope anew.

\section{Whole components, derived endpoints and determinant relations}
The full node lattice category has precisely the three indecomposables $X,\Lambda,Y$ and arrows $X\to\Lambda\to Y$, $Y\to\Lambda\to X$. Its stable category has only $X,Y$, no irreducible arrows, and suspension exchanging them; its AR triangles have zero middle term. For the polynomial node at a smooth support $x=\lambda+t$, $y=0$, $\lambda\ne0$, the modules $T_r=k[t]/t^r$ form an ambient homogeneous tube:
\[
 0\longrightarrow T_r\xrightarrow{(\pi,t)}T_{r-1}\oplus T_{r+1}
 \xrightarrow{(t,-\pi)}T_r\longrightarrow0,\quad T_0=0,\quad\tau T_r=T_r.
\]
The shift changes the derived placement, not the quasi-length. These nonzero-central-support objects do not extend to continuous finite-length modules of the arrow-adic completion. At its singular point even the candidate pure-branch tube fails ambient almost-splitness: $T_2^y\to k$ cannot factor through $T_2^x\to k$, since every intervening map lands in the $x$-socle, killed by the last map.

The nine HLS strings above sit in an entire intrinsic quotient component $T_{p,q}$, $(p,q)\in\Z^2$, with arrows increasing one coordinate and $\tau T_{p,q}=T_{p-1,q-1}$ \cite{HuangLinSu}. Its central object $E=M(ab)$ has infinite projective dimension: restriction to the first corner gives the dual-number simple with its infinite loop resolution, while restricted projectives are free or zero. Its module AR sequence therefore yields a distinguished derived triangle but not a derived AR triangle ending at $E$. Ringel's four-quarter wind-wheel tile \cite{Ringel}, with infinite-word diagonals and mixed center, is a different category's quilt; it is not silently identified with that HLS component or the Belyi surface.

There is no general equation ``AR mesh equals polynomial semi-invariant syzygy.'' The intrinsic $k(1\to2)$ quotient of the infinite Chebyshev algebra has AR row $0\to S_2\to E\to S_1\to0$ and a free presentation determinant $z$ on dimension $(1,1)$, so its semi-invariant ring $k[z]$ has no generator relation. A genuine framed elliptic-square stratum has $Z=[x_1,x_2,y_1,y_2]\in\operatorname{Mat}_{2\times4}$ with six minors and relation
\[
 p_{12}p_{34}-p_{13}p_{24}+p_{14}p_{23}=0.
\]
These minors have the same determinant character for the target $\mathrm{GL}_2$ with framed source; they are not all semi-invariants for the full unframed vertex group. In contrast, the nodal lattice AR row gives a genuine relative determinant identity over $R$: multiplication by $a+bx$ on the basis $(y,1)$ has matrix
$\bigl(\begin{smallmatrix}a&-b\\0&a+ub\end{smallmatrix}\bigr)$ and characteristic factor $(1-Ta)(1-T(a+ub))$. The endpoints have infinite order-projective dimension but are $R$-perfect.

Finally, the three-vertex theta algebra has no unital surjection onto a six-vertex radical-square-zero circular algebra: its semisimple image has at most three primitive idempotent images. Duplicating face vertex spaces instead gives an exact faithful unfolding which is not full; it does not inherit quotient or AR-transfer claims. Zwara Hom-control means that the Hom-dimension defect is bilinear on $K_0$, not merely a covering direct-sum formula \cite{ZwaraCorrected,RiedtmannZwaraCorrected}. The book supplies explicit fixed-dimension counterexamples, corrected circular-complex incidence resolutions, complete drawn component windows and all category hypotheses. These additions do not supply the arithmetic determinant comparison.



\chapter{Verification and research frontier}
The two nodal/invariant Lean files construct the node as an actual subring of $k[[X]]\times k[[X]]$ with equal constants, prove its annihilators, regular common parameter, augmentation, alternating projective resolution and a bijective pair of power-series coordinates. A second file proves CK1 collapse under the stated permutation/coverage hypotheses, determinant base change, actual $2\times2$ symmetric-square functoriality and relation preservation, a Schur character identity and the fixed-parameter lattice failure. These thirty printed theorem types have only standard Lean foundations. They do not formalize general dessins, abstract Ext/AR/Schur--Weyl, C* exactness, Fredholm determinants, arithmetic descent or GRH. Earlier formal layers remain preserved with their literal scopes.

Exact-rational diagnostics independently check the displayed dessins, path identities, Schur and Young examples and the nine-object grid. These finite checks supplement the proofs, rather than prove infinite spectral or arithmetic assertions. The next arithmetic target is an exact descent/Frobenius-compatible construction on the labelled order/Fock data, with a determinant-preserving positive comparison. The local Schur kernel and the simple degree spectrum alone do not supply it. All source clones and earlier research are retained; the earlier general focused PDFs are archived at the preceding commit.

Ten further declarations in \texttt{DessinPullbackComponents.lean} check actual matrix-unit identities, nonmultiplicativity of raw diagonal extraction, impossibility of a unital two-state matrix residue into a nonzero commutative ring, multiplicativity on diagonal fibers, and genuine equal-image pullback subring/projection components. The analytic suite then contained187printed declarations at unchanged pins. The circle-function C* model, functional calculus, generation, graph uniqueness and K-theory remain prose/external-theorem arguments, not hidden assumptions of these component checks.

Sixteen additional declarations in \texttt{ThetaJacobi.lean} check genuine finite-field MulChar Jacobi conjugation, norm and rescaling with explicit nontriviality hypotheses; cubic--quadratic noncancellation; the cyclotomic multiplication determinant and positive Gram identity; and actual theta cellular boundary, homology-inclusion, restriction and rational chain-map matrices. The full analytic suite now contains203declarations. The point-count proof, deck/Kummer comparison, cohomological trace interpretation and global arithmetic comparison are not thereby fully formalized.

The final three \texttt{ThetaJacobi.lean} declarations verify the actual conjugation matrix identities and the real two-dimensional inert-prime determinant obstruction. They do not assume any unknown zero location or exclude complex/enlarged spectral realizations.


\texttt{ThetaProjectiveMultiplicity.lean} adds five declarations. Arbitrary upper-block matrices with six identical diagonal blocks have the sixth-power determinant and characteristic-factor formulas. A literal power-series constant-coefficient map specializes matrix determinants, for each scalar parameter; the coefficientwise polynomial/order interface remains prose. A nonzero quadratic cannot be a rational-function sixth power, proved by polynomial degrees. Three genuine power series with product $X$ cannot have the cyclically equivalent residue-vanishing conditions forced by invertible rank-one arrow intertwiners. The derivation of the path blocks from the full theta order, projective decomposition, derived branch functor and deck intertwiners is not fully formalized. The complete analytic suite now has208printed declarations; no hidden arithmetic positivity assumption occurs.


\texttt{ThetaTameDescent.lean} adds twelve printed component declarations (and one auxiliary polynomial lemma). They use genuine formal-power-series rescaling and its coefficient operator, proving the coefficient action, invariant-coefficient vanishing and the tame operator's cube. Actual three-state matrices satisfy the clock/rotation commutator, comparison determinant and closed-fiber equations, weighted-cycle cube and full characteristic polynomial, noncommutation of a nonconstant rotation polynomial, trivial whole-column determinant actions, and the same-operator comparison for $vC$. The classification of all fixed series by residues, full order action and hereditary invariant algebra, stack interpretation, valuation/rank theorem, general coefficient-functor filtration and GRH are not thereby fully formalized. No unknown positivity premise is assumed.

\part{Proofs and calculations}

\chapter{Expanded dessin, resolution and derived calculations}
\section{The cycle order in coordinates}
Number a permitted cycle in travel order $0,\ldots,\ell-1$. Map an arrow $i\to i+1$ to $E_{i,i+1}$ for $i<\ell-1$, and the final arrow to $u_c E_{\ell-1,0}$. A path from $i$ to $j$ then gives $u_c^rE_{ij}$, where $r$ counts cut crossings. These are linearly independent path normal forms, and every positive-length permitted path has exactly one such expression on exactly one branch. The full-cycle sum in a branch is $u_c I$. At a paired edge its two diagonal occurrences come from a single original idempotent; no positive-length coefficients are identified. This proves the order description including self-pairings within one cycle. Although the suborder is generically a product of full matrix algebras, it is not their unglued product at the special fiber.

For the interval this becomes the literal fiber product
\[
 \Lambda=\{(f,g)\in k[[X]]^2:f(0)=g(0)\},\qquad
 x=(X,0),\ y=(0,X),\ u=(X,X).
\]
Every element has a unique expression $a(u)+b(u)x$: take $a=g$ and $b=(f-g)/X$. Thus it is $R$-free on $1,x$ with $x^2=ux$ and $y=u-x$. The annihilator of $x$ consists of pairs $(0,g)$ with $g(0)=0$, exactly $y\Lambda$; the other branch is symmetric. Multiplication by $u$ is injective, coefficientwise. These are actual formal power-series statements, not conclusions from finite truncations.

\section{The two-arm complex and its calculations}
For an arrow $a$ write $f(a)$ for its unique forbidden successor. The summand in $P_n(i)$ indexed by the initial arrow $a$ has endpoint $t(f^{n-1}(a))$. Its map multiplies by that last arrow. Thus the resolution can be computed by iterating a finite permutation, while the projectives themselves remain infinite-dimensional. The Hom complex against a module $N$ is the direct sum of the endpoint spaces with differentials given by the same arrow matrices. Tensor against a left module instead uses the same path entries acting on the left. No rank or exactness claim for an arbitrary $N$ is inferred by counting endpoints alone.

For the nodal residue field the maps are particularly explicit:
\[
 \cdots\xrightarrow{\operatorname{diag}(x,y)}\Lambda^2
 \xrightarrow{\operatorname{diag}(y,x)}\Lambda^2
 \xrightarrow{(x\ \ y)}\Lambda\longrightarrow k\longrightarrow0.
\]
If $xz_1+yz_2=0$, its first branch says the first branch of $z_1$ is zero, and its second branch says the second branch of $z_2$ is zero. Hence $z_1=yw_1$, $z_2=xw_2$. The same annihilator calculation proves the alternating kernels in all degrees. Projectivity is genuine: the terms are regular modules or their finite direct sums. It is not a finite free resolution, despite the period of its positive tail.

For the elliptic theta the quiver vertices are $v_0=\{0,3\}$, $v_1=\{1,4\}$, $v_2=\{2,5\}$. The arrows $a_0,a_3$ go $v_0\to v_2$, $a_2,a_5$ go $v_2\to v_1$, and $a_1,a_4$ go $v_1\to v_0$. The two forbidden arms from $v_0$ are
\[
 a_0,a_5,a_1,a_3,a_2,a_4,\ldots,\qquad
 a_3,a_2,a_4,a_0,a_5,a_1,\ldots.
\]
Their top pairs are $(v_2,v_2),(v_1,v_1),(v_0,v_0)$, repeated. The three steps exchange the arms. Their distinct top pairs exclude shorter unordered positive periods. The spherical theta instead has forbidden two-cycles and successive top pairs $(v_2,v_1),(v_0,v_0)$, giving period two. Thus even identical hereditary overorders need not give identical resolutions.

Quadratic Koszul duality here has a directly checkable local meaning: relations complement the two possible successors. The dual dessin replaces the graph rotation by $\phi^{-1}$. In the Ext algebra against $A/J$, concatenating the forbidden-word basis gives the product when the join is forbidden; a permitted join gives zero. One can see this by lifting a length-$r$ basis class to the resolution: discard the prescribed initial $r$ letters and identify the remaining summand with its suffix resolution, taking zero on other prefixes. These prefix-discard comparison maps commute with the suffix differentials, with the standard shift signs when viewed as shifted complexes. The concatenated indexing word exists exactly when its join is forbidden; evaluation of the two basis classes on that word gives one. If the join is permitted there is no such indexing word and the product is zero. Composition therefore concatenates indexing words, with the chosen right-module/Yoneda opposite convention. This is a graded algebra statement; it neither makes the unbounded resolution perfect nor constructs an equivalence of all formal-completion derived categories.

For an arbitrary bounded projective complex $P^\bullet$, its determinant line over a commutative coefficient ring is $\bigotimes_n(\det P^n)^{(-1)^n}$, with determinant duals for negative exponents. A chain automorphism acts on this line; a quasi-isomorphism gives the corresponding determinant-functor isomorphism. For a general endomorphism $F$, the expression for $1-zF$ uses the locus where it is invertible and gives a rational determinant-line action/trivialization, not an everywhere invertible line map at singular parameters. Additivity is compatible with exact triangles under the usual determinant-functor conventions \cite{KnudsenMumford,Breuning}. In our order context an $R$-linear determinant uses the finite free underlying $R$-modules; it must specify a separate endomorphism. Determinants of a differential presentation, alternating determinants of a chain action, regularized analytic determinants and the Cartan path matrix are four distinct constructions. Their equality requires a proof, not their shared name.

The geometric model literature \cite{GeometricGentle,GeometricResolutions,GeometricKoszul} is useful for identifying which crossings and gradings would have to be represented by a functor. To apply it here one must specify the category of polynomial strings, completed lattices or finite-length quotients; construct its surface/curve objects; and prove that degree-$r$ intersections compute the Hom group into shift $r$, with local systems and boundary exceptions included. We have supplied explicit walks and maps in the examples, not claimed this general completed-order theorem. On a finite word, a map can already be computed rigorously without such a theorem: assign a coefficient matrix between equal-vertex occurrences and solve $F_{t(a)}N(a)=N'(a)F_{s(a)}$ for every arrow. The diamond calculations below implement exactly this system.


\chapter{Semi-invariants, tableaux and explicit intertwiners}
\section{Derived information retained by a presentation}
Let $K=\operatorname{im}f$. The injection $\operatorname{Hom}(K,N)\to\operatorname{Hom}(P_1,N)$ descends to the Ext injection used in Theorem~\ref{thm:semi}; its kernel modulo $\operatorname{im}\operatorname{Hom}(P_0,N)$ is zero because the map $P_1\to K$ is onto. The missing part of the cokernel detects maps on $P_1$ not descending through $K$. This explains exactly why an injective resolution map is needed for its full Euler interpretation.

In the node example the resolution of $M=\Lambda/(x)$ has maps $x,y,x,y,\ldots$. On $N=k[y]/(y^2)$ the Hom cochain maps are $0,y,0,y,\ldots$. Thus $\operatorname{Hom}(M,N)=N$ and $\operatorname{Ext}^1(M,N)=\ker y$ have dimensions two and one; the cokernel of the first zero map is two-dimensional. In the elliptic-square example the first four Hom maps are $1,0,0,0$, periodically. Hence degrees zero and one vanish, degrees two and three are $k$, and degree four vanishes. This also makes clear why a nonzero scalar semi-invariant does not give a derived orthogonality theorem.

With coordinate action $(gF)(N)=F(g^{-1}N)$ the Schofield weight is $a-b$, the negative of the point-evaluation weight. For a single arrow it is the same elementary distinction as the action on $\operatorname{Hom}(V,W)$ versus on its coordinate functions $V\otimes W^*$. The Cauchy decomposition for coordinate functions is
\[
 \operatorname{Sym}^m(V\otimes W^*)\cong
 \bigoplus_{\lambda\vdash m}S_\lambda(V)\otimes S_\lambda(W)^*.
\]
A relation quotient can kill these components; its ambient multiplicities do not descend unchanged. Acyclic-quiver semi-invariant/LR results such as \cite{DerksenWeyman} retain their acyclicity assumptions. The explicit two-arrow slice below is acyclic, whereas the whole theta quiver is not.

\section{One Young operator and one LR decomposition in coordinates}
For $\lambda=(2,1)$ number the first row $1,2$ and second row $3$. Then
\[
 Y=(I+P_{12})(I-P_{13}),\quad Y^2=3Y,\quad
 \dim\operatorname{im}(Y\mid(k^2)^{\otimes3})=2.
\]
The hook product is three; its Specht dimension is $3!/3=2$, also two here for a different reason. The normalized projector extracts one $S_{(2,1)}$ copy. The trace $\tr((Y/3)M^{\otimes3})$ is therefore the Schur character, not twice that character. This fixes the common row/column normalization ambiguity.

Write $t=\tr M$, $\Delta=\det M$ for a two-dimensional endomorphism. Its complete characters obey $h_0=1$, $h_1=t$, $h_m=th_{m-1}-\Delta h_{m-2}$, and
\[
 \chi_{(a,b)}(M)=\Delta^b h_{a-b}(M),\quad a\ge b\ge0.
\]
In particular $h_2=t^2-\Delta$, $h_3=t^3-2t\Delta$, and $\chi_{(2,1)}=t\Delta$. Thus $t h_2=h_3+\chi_{(2,1)}$ is the LR identity $V\otimes\operatorname{Sym}^2V\cong\operatorname{Sym}^3V\oplus S_{(2,1)}V$, with $S_{(2,1)}V\cong\det V\otimes V$.

An explicit natural splitting avoids relying just on characters. In characteristic zero multiplication $m:V\otimes\operatorname{Sym}^2V\to\operatorname{Sym}^3V$ has natural section
\[
 s(v_1v_2v_3)=\tfrac13\sum_{i=1}^3v_i\otimes\prod_{j\ne i}v_j.
\]
The map $j:\bigwedge^2V\otimes V\to V\otimes\operatorname{Sym}^2V$ sends $(v\wedge w)\otimes z$ to $v\otimes wz-w\otimes vz$. In dimension two it is injective: for a basis $e_1,e_2$ its images at $z=e_1,e_2$ have distinct independent coordinates. They lie in $\ker m$, whose dimension is $6-4=2$. Hence $(s,j)$ is the required isomorphism. Tensorial definitions make both maps commute with every linear endomorphism, including singular arrow maps. This is a genuine LR intertwiner on our dessin representation stratum.

For a concrete theta representation take
\[
 X_1=\begin{pmatrix}1&1\\0&1\end{pmatrix},\quad
 X_2=\begin{pmatrix}2&0\\0&1\end{pmatrix},\quad
 X_3=\begin{pmatrix}1&0\\1&1\end{pmatrix}.
\]
Then $M=X_3X_2X_1=\left(\begin{smallmatrix}2&2\\2&3\end{smallmatrix}\right)$, $t=5$, $\Delta=2$, $h_2=23$, $h_3=105$ and $\chi_{(2,1)}=10$. The identity is $115=105+10$. Since the other permitted branch is zero, every forbidden product is zero; this is an actual representation of the infinite elliptic dessin algebra.

In the invertible-arrow open stratum one can gauge $X_1,X_2$ to the identity, leaving conjugation of $M$. The invariant-ring proof in the core works also at singular arrows by cycle-trace generation, not by assuming that this gauge choice extends there. The degree-one monodromy coefficients have total arrow degree three; this is why the Hilbert-series denominator uses $3j$, not $j$.

For the Cauchy kernel choose the Schur subspace inside the tensor power with its inherited Hilbert norm. $M^{\otimes m}$ preserves it and has norm at most $\|M\|^m$. This proves the trace bound for nonnormal matrices as well as diagonal ones. Formal Cauchy identities extend by polynomial continuity in matrix entries; the stated absolute bound then justifies analytic summation. Gram positivity does not assert entrywise positivity or positivity of an unrelated arithmetic form. For three exact test matrices $M$, $\left(\begin{smallmatrix}1&1\\0&2\end{smallmatrix}\right)$, $\left(\begin{smallmatrix}0&1\\-1&1\end{smallmatrix}\right)$, the choice $r=1/100$ satisfies $rR_0^2\le21/100<1$ using Frobenius norms, and all seven principal minors of the rational kernel matrix are nonnegative.

\section{Two precise functorial limitations}
At any actual vertex simple $S$ of $A$, the split exact sequence $0\to S\to S\oplus S\to S\to0$ uses first inclusion and second projection. Termwise $\operatorname{Sym}^2$ has dimensions $1,3,1$ at that vertex; its middle kernel contains the mixed monomial as well as the first square, but its first image contains only the first square. Middle homology is $S$. Thus even split exactness is not preserved. Genuine derived Schur constructions instead use mixed tensor degrees, signed total differentials and suitable flat/projective replacements \cite{SchurComplexes,DerivedSchur}; their membership in a particular order category must still be checked.

The $R$-lattice $M=R e_x\oplus R e_y$ with $x=\operatorname{diag}(u,0)$, $y=\operatorname{diag}(0,u)$ has $x+y=uI$. On $\operatorname{Sym}^2_RM$ the images are $\operatorname{diag}(u^2,0,0)$ and $\operatorname{diag}(0,0,u^2)$. Their sum kills the mixed vector, hence equals neither $uI$ nor $u^2I$. This is a specific failure to retain the original central action, not a failure of the polynomial monomial representation functor. Even a common ramified base substitution does not repair this mixed vector.

Semi-invariants retain useful moduli on nilpotent modules. In the theta quiver keep only $a_0,a_3:v_0\to v_2$, put their scalar values $(x,y)=(1,\lambda)$ and all other arrows zero. The $v_1$ simple splits off. Every length-two path acts by zero, so these are continuous completed modules. The two scalars have the same weight $g_{v_2}/g_{v_0}$, and $y/x$ is a rational invariant on $x\ne0$, giving the projective pencil $[x:y]$. It is not a regular invariant on the whole affine locus, whose orbit closures meet zero.

For the interval quotient take $x=J_2$, $y=\lambda J_2$. With its top and socle specified, the separated Kronecker coordinates are again $(1,\lambda)$; the relation to the same-weight ratio is exact. The pp formula in the core is stronger than merely noticing a ratio: on $M_\lambda$ its solution set is $\operatorname{im}x$, whereas on $M_\mu$, $\mu\ne\lambda$, the equation $(\mu-\lambda)xv=0$ forces $z=xv=0$. Different parameters are therefore distinguishable by actual positive-primitive conditions even though every cycle trace is zero.


\chapter{Full lattice and diamond diagrams}
\section{The nodal exact diagram}
The normalization proof of Theorem~\ref{thm:nodear} uses $BM$ inside the generic module. Replacing it without justification by $B\otimes_\Lambda M$ could introduce torsion and would not prove the same statement. Once the subdirect residue subspace is decomposed, the branch basis changes lift over complete DVRs; taking inverse images modulo the conductor gives the direct decomposition, not just a decomposition of residues.

Here is the promised diagram, with the terminal short exact column for every position:
\[
\begin{tikzcd}[column sep=large,row sep=large]
0\arrow[r]&Y\arrow[r,hook]\arrow[d,"u"]&\Lambda\arrow[r,"x"]\arrow[d,"u"]&X\arrow[r]\arrow[d,"u"]&0\\
0\arrow[r]&Y\arrow[r,hook]\arrow[d]&\Lambda\arrow[r,"x"]\arrow[d]&X\arrow[r]\arrow[d]&0\\
0\arrow[r]&Y/uY\arrow[r,hook]&\Lambda/u\Lambda\arrow[r]&X/uX\arrow[r]&0
\end{tikzcd}
\]
Each vertical column is completed by zero above and below. In the bottom row identify $Y/uY=k$ via generator $y$, $X/uX=k$ via $x$, and $\Lambda/u\Lambda=k[\epsilon]/(\epsilon^2)$ with $x=\epsilon$, $y=-\epsilon$. Its first map is $1\mapsto-\epsilon$, its second the augmentation. It is nonsplit as a dual-number sequence. This explicit order reduction, not an unrelated finite algebra, explains its appearance.

For a further Ext calculation, resolve $X\cong\Lambda/(y)$ with successive maps $y,x,\ldots$. On $Y$ the first is multiplication by $y$ and the second zero, so $\operatorname{Ext}^1(X,Y)=Y/yY=k$; on $X$ the first is zero and the second injective, so $\operatorname{Ext}^1(X,X)=0$. Interchanging branches gives the other two cases. The two nonsplit lattice rows realize the nonzero cross classes. Their Hom spaces are finite over $R$, generally infinite over $k$.

\section{Chebyshev geometry and the exact quotient}
The identity $1-T_n(z)^2=(1-z^2)U_{n-1}(z)^2$ follows from $z=\cos\theta$ and the sine/cosine multiple-angle formulas, hence as a polynomial identity. $T_n'=nU_{n-1}$. Therefore $\beta_n$ has its only zero at infinity, of order $2n$; its $n$ points over one are double since $\beta_n-1=T_n^2/(1-T_n^2)$; and its poles are simple at $\pm1$ and double at the $n-1$ interior roots of $U_{n-1}$. The derivative has no additional critical values. The total ramification is $(2n-1)+n+(n-1)=4n-2$, as required for a degree-$2n$ rational Belyi map. The dual graph has $n+1$ vertices of degrees $1,1,2,\ldots,2$, and $n$ edges; it is a connected path. Its planar ribbon embedding is unique, so this identifies the particular nested-loop dessin, not only its passport.

In the quotient retaining orientation $\Omega$, any surviving path traverses each chain edge at most once and either endpoint loop at most once. Its length is at most $n+1$; thus the killed-arrow ideal contains $J^{n+2}$. The quotient of $A_n$ and of its completion is the same $H_n(\Omega)$. For $n=3$ with forward chain arrows, its path counts in lengths zero through four are $3,4,3,2,1$, total thirteen. These counts describe a quotient representation category of the infinite order, not a change of primary object.

Inflation is full because the two algebras have exactly the same actions on these annihilated objects. It preserves internal exact sequences. It does not preserve arbitrary ambient extension closure: the two adjacent simples are annihilated by the killed reverse arrow, but the two-dimensional extension having only that arrow nonzero is not. Likewise endpoint-local freeness over $k[l]/l^2$ or $k[r]/r^2$ in \cite{HuangLinSu} is not $R$-freeness. To transfer an almost-split property to a larger category one must prove the requisite approximation/factorization statements; locally finite quiver criteria \cite{BautistaLiuPaquette,Paquette,LiuNgPaquette} do not apply just because each vertex has two arrows. The permitted cycles violate interval/Hom finiteness.

\section{Nine actual strings and all maps}
Take $n=3$ with $1\xrightarrow{a}2\xrightarrow{b}3$, loops $l,r$ and $l^2=r^2=0$ in the quotient. Here $b$ denotes the retained forward $a_2$, not a reverse arrow of the original doubled chain. The following words are in travel order. A superscript $-$ means a formal inverse in a word, not an inverse algebra element.
\begin{center}\small
\begin{tabular}{llll}
\toprule
Object & word & start & dimension vector\\\midrule
$A$&$l^- a b r b^- a^- l^-$&1&$(4,2,2)$\\
$B$&$l^- a b$&1&$(2,1,1)$\\
$C$&$l^- a b r^-$&1&$(2,1,2)$\\
$D$&$a b r b^- a^- l^-$&1&$(3,2,2)$\\
$E$&$a b$&1&$(1,1,1)$\\
$F$&$a b r^-$&1&$(1,1,2)$\\
$G$&$r^- b^- a^- l a b r b^- a^- l^-$&3&$(4,3,4)$\\
$H$&$r^- b^- a^- l a b$&3&$(2,2,3)$\\
$I$&$r^- b^- a^- l a b r^-$&3&$(2,2,4)$\\\bottomrule
\end{tabular}
\end{center}
A basis vector is assigned to each visited position $0,\ldots,\ell(w)$ at its vertex. A direct letter sends its preceding vector to its following vector; an inverse letter sends the following vector to the preceding vector. Sum these elementary maps for each arrow and put all unmentioned entries zero. This defines every matrix, including repeated loops, and verifies the relations. Set the killed reverse arrows to zero to obtain actual $\Lambda_3$-modules.

Their signed grid is
\[
\begin{tikzcd}[column sep=large,row sep=large]
 A\arrow[r,"p"]\arrow[d,"-p"']&B\arrow[r,"i"]\arrow[d,"p"]&C\arrow[d,"-p"]\\
 D\arrow[r,"p"]\arrow[d,"-i"']&E\arrow[r,"i"]\arrow[d,"i"]&F\arrow[d,"-i"]\\
 G\arrow[r,"p"]&H\arrow[r,"i"]&I.
\end{tikzcd}
\]
Horizontal projections keep the initial word positions and kill the remaining positions. Horizontal inclusions append the final $r^-$, keeping old positions. Upper vertical projections delete the initial $l^-$: old position $j\ge1$ maps to position $j-1$, and old position zero maps to zero. Lower vertical inclusions prepend $r^-b^-a^-l$: position $j$ maps to $j+4$. Multiply these maps by the indicated signs. For every arrow these formulas commute with the string action, including the two positions adjacent to a cut or join; those are exactly where the loop inverse and chosen orientation matter.

Every small square has two length-two composites whose sum is zero. The exact sequences are
\[
\begin{split}
0&\to A\to B\oplus D\to E\to0,\qquad
0\to B\to C\oplus E\to F\to0,\\
0&\to D\to E\oplus G\to H\to0,\qquad
0\to E\to F\oplus H\to I\to0.
\end{split}
\]
For example the maps out of $A$ retain its initial positions $0,1,2,3$ in $B$ and delete its first position in $D$, so their joint kernel is zero. The maps to $E$ are onto, and the dimension vectors give $\dim A+\dim E=\dim B+\dim D$ vertexwise. The zero composite plus this equality proves exactness. The other three use precisely the same keep/delete or append/prepend rules; their vertexwise dimension equalities and jointly injective/surjective maps prove exactness. The shipped exact checker constructs all matrices and verifies these assertions independently over $\mathbb Q$.

The stronger almost-split property is imported by identifying $E=w=ab$ as the minimal string of type $(2,2)$ and these exact hook/cohook words with \cite[Proposition3.12]{HuangLinSu}, using that paper's Propositions2.6/2.9 and Lemma3.10. Opposite-module conventions are reconciled by reversing travel/composition conventions; the closed doubled quiver is isomorphic to its opposite by exchanging the two arrows at each chain edge. Thus intrinsically $\tau E=A$, $\tau F=B$, $\tau H=D$, $\tau I=E$. The rows and columns of this grid are \emph{not} generally short exact, unlike the nodal diagram. Two small cautions in that source are not needed here: a tube-rank label in Example3.9(3) is inconsistent with its listed mouth, and a general homogeneous-tube statement must not be read as a converse classifying every such object as a band. Neither is used to import these four meshes.

\section{The pure-injective quotient and what its dimension means}
For $B_0=\Lambda/J^2=k\langle x,y\rangle/(x,y)^2$, words alternate direct and inverse letters because every direct/direct and inverse/inverse adjacent pair is forbidden. Up to rotation and inverse its only primitive band is $xy^{-1}$; the bridge quiver has two oriented band vertices with no connecting nontrivial bridge. The longest path has length zero. The domestic string results \cite{PuninskiPrest,LakingPrestPuninski} therefore give relative dimension $0+2=2$, not an automatic dimension three for every one-domestic algebra. Finite-dimensional points have relative CB rank zero; the appropriate Prüfer/adic and one-sided periodic string points have rank one; the unique generic point has rank two. These statements concern the quotient category with its inherited pp structure.

Indeed a pp formula is a finite system of linear equations with existentially quantified unknowns. On a module annihilated by $J^2$ every coefficient reduces to $B_0$, so purity and pure-injectivity are preserved in both directions under inflation. The condition $MJ^2=0$ is a definable closed condition, giving the closed Ziegler subspace. Its relative ranks need not be ranks in the entire $\Lambda$ spectrum. This subcategory is not extension closed: $\Lambda/(x^3,y)$ gives an extension of objects killed by $J^2$ that is not itself killed by $J^2$. Since $u^2$ annihilates every object of the subcategory, it contains no nonzero $R$-free lattice. Ringel's representation-theoretic node construction is also not a Morita equivalence identifying the geometric nodal order with the Kronecker algebra. Infinite-dimensional polynomial string classifications \cite{CrawleyBoeveyStrings} retain finitely-controlled or pointwise-artinian assumptions; an arbitrary complete lattice does not automatically satisfy them as an algebraic $A$-module.


\chapter{The operator extension in detail}
\section{A corrected cycle lift and the integer boundary map}
In the permitted-state graph each length-$\ell$ cycle has a full matrix-corner Laurent/circle coordinate. Its chosen $K_1$ generator is the full cycle at one state plus the complementary identity. In Fock space $L_c^*L_c=Q_a+v_{s(a)}$. Thus the tempting lift $L_c+(1-Q_a)$ has overlapping initial pieces and is not an isometry. The corrected lift in Theorem~\ref{thm:fock} first restricts to $Q_a$:
\[
 W=L_cQ_a+(1-Q_a),\quad W^*W=I,\quad
 I-WW^*=Q_a-L_cQ_aL_c^*.
\]
A path starting with $a$ and length greater than $\ell$ has uniquely the full prefix $c$ followed by another path starting with $a$. The missing range consists exactly of lengths $1,\ldots,\ell$. Its terminal multiset is the original graph's incident-edge multiset at that cycle vertex. These finite defect vectors all lie in terminal compact ideals; their classes give the claimed unsigned incidence entries. The sign convention affects only the overall sign of the boundary map.

Here is the Smith calculation without spectral assumptions. Choose a spanning tree of the connected underlying graph and one root. By changing signs on one bipartition of the tree, its $E\times V$ incidence submatrix becomes the ordinary oriented tree incidence, with $V-1$ unit pivots. After eliminating these pivots, a remaining column represents the common root coordinate, and each nontree edge evaluates on it as zero when its fundamental cycle is even, or $\pm2$ when odd. A loop is the latter case. If all are even the graph is bipartite and this column is zero; otherwise its entries have gcd two, giving the final pivot two. Remaining rows are free. Consequently the kernel/cokernel, then the six-term sequence, give the $K$ groups stated in the core. No theorem reconstructing a cyclic embedding from those groups is asserted.

The compact ideal is block diagonal, not $\mathcal K(\mathcal H)$: prefixes never move terminal vertices, so cross-terminal rank-one maps do not belong to the displayed ideal. This is why the arithmetic insertion must choose one terminal sector. A graph/Cuntz boundary quotient kills the vacua and hence this ideal, so that insertion need not survive passage to the quotient.

\section{Spectral domain and arithmetic trace}
A vector belongs to the adjoint domain of the diagonal $N$ precisely when its coefficientwise product by the real sequence $|p|$ is square summable; this proves self-adjointness on the given domain. The truncations by length prove finite supports are a core. Resolvent coefficients $(|p|-i)^{-1}$ tend to zero and each length space is finite dimensional, giving compactness. Prefix and deletion shift the length by at most one, so preserve the domain and prove the commutator formulas there. Leibniz then handles the dense polynomial *-algebra.

Every terminal sector has one vacuum and two paths at each length $n\ge1$, even when some graph endpoints coincide: distinct formal paths remain distinct basis vectors. Thus restriction to a sector still has a simple vacuum, but it need not be a faithful representation of the whole concrete algebra. The trace of $e^{-tN^2}$ is $E+2E\sum_{n\ge1}e^{-tn^2}$, distinct from the semigroup trace of $N$ in the core. The trace identity with $\chi(N)$ is only on $\Rea s>1$, where the product is trace class. Analytic continuation of the scalar $L$-function does not make that same operator product trace class on $\Rea s>0$.


\chapter{Abel determinant estimates and the arithmetic program}
For completeness the inserted family has a full elementary convergence proof. Since complete nonprincipal character periods sum to zero, $|A_n|\le B_0\le q_\chi$. On $2^k\le n<2^{k+1}$,
\[
 \sum |u_n|^2\le\frac{B_0^2}{(k+1)^2},\qquad
 \sum |v_n(s)|^2\le M^2(k+1)^2r^k,\quad r=2^{-2\sigma_0}<1.
\]
The second follows from $n^{-s}-(n+1)^{-s}=s\int_n^{n+1}x^{-s-1}\,dx$. Define
\[
 Q_J(r)=r^J\left(\frac{(J+1)^2}{1-r}
 +\frac{2(J+1)r}{(1-r)^2}+\frac{r(1+r)}{(1-r)^3}\right).
\]
Expanding $(J+j+1)^2$ in the three geometric moment sums proves $Q_J=\sum_{k\ge J}(k+1)^2r^k\to0$. Locally uniform $\ell^2$ convergence of the holomorphic finite-coordinate vectors proves Banach holomorphy. The linear functional $z\mapsto\sum v_nz_n$ has norm $\|v\|$; this is the complex-linear coefficient pairing, not a conjugate-linear dependence on $s$.

Write $\mathcal I(v)(z)=\sum_n v_nz_n$ for the coefficient functional, distinct from the arrow ideal $J$. For $N_J=2^J-1$, $J\ge1$, let $K_N=u\otimes(\mathcal I(P_Nv)-e_1^*)$. Then
\[
 \|K-K_{N_J}\|_1\le\sqrt2 B_0M\sqrt{Q_J(r)}.
\]
Finite summation by parts gives
\[
 T_N=\det_F(I+K_N)
 =\sum_{n\le N}A_n(n^{-s}-(n+1)^{-s})
 =S_N-A_N(N+1)^{-s}.
\]
Adding the last endpoint times $e_1^*$ to the functional recovers $S_N$ exactly because $u_1=1$. Its trace-norm error is at most
\[
 \sqrt2 B_0\bigl(M\sqrt{Q_J(r)}+B_0(N_J+1)^{-\sigma_0}\bigr).
\]
For arbitrary cutoffs use the last whole dyadic block for an enclosing tail. Galerkin compressions also converge: $\|u-P_{N_J}u\|\le B_0/\sqrt J$. With $C_K=1+M\sqrt{(1+r)/(1-r)^3}$,
\[
 \|K-P_{N_J}KP_{N_J}\|_1
 \le\frac{B_0 C_K}{\sqrt J}+\sqrt2 B_0M\sqrt{Q_J(r)}.
\]
The raw endpoint adds the same endpoint bound as before. Isometric insertion by $U$ preserves nonzero singular values, every norm estimate, and the Fredholm determinant. Everything is in one compact ideal of the actual infinite dessin Fock algebra.

For primitive parity $a$ use $g(s)=(q_\chi/\pi)^{(s+a)/2}\Gamma((s+a)/2)$, holomorphic nonzero on $\Rea s>0$. The completed functional is $g(s)\mathcal I(v)-e_1^*$, and the completed raw functional is $g(s)(\mathcal I(P_Nv)+A_N(N+1)^{-s}e_1^*)-e_1^*$. Their determinants are exactly $gL$ and $gS_N$. Scaling the entire old perturbation would instead give $1+g(L-1)$, the wrong normalization. Trace-class background is standard \cite{Bornemann}, while these special rank-one identities follow directly from $\det_F(I+u\otimes\varphi)=1+\varphi(u)$.

The metric obstruction is likewise elementary. With inner product conjugate-linear in its first variable, symmetry of $T=u\otimes\varphi$ implies
\[
 \varphi(z)\langle u,u\rangle=\overline{\varphi(u)}\langle u,z\rangle.
\]
A nonzero $\varphi$ cannot have $\varphi(u)=0$; thus all symmetric maps with that range have proportional functionals. For $n\ge2$ the coefficient ratios at $\sigma+it$ and $\sigma$ satisfy
\[
 \frac{v_n(\sigma+it)}{v_n(\sigma)}
 =\frac{\sigma+it}{\sigma}\,n^{-it}(1+O_{\sigma,t}(1/n)).
\]
Integers nearest $\exp(2\pi k/|t|)$ and $\exp((2\pi k+\pi)/|t|)$ give opposite nonzero subsequential limits. Hence no common positive metric symmetrizes this family for all critical heights. Restricting any proposed metric on the inserted family to $U\ell^2$ gives the same contradiction. An independently constructed positive comparison is still a possible different route.

\section{Precisely how classical and recent sources enter}
Over a finite field the Grothendieck trace-formula architecture writes a zeta function as alternating determinants on finite-dimensional cohomology; Deligne's purity supplies the weights of eigenvalues \cite{SGAFourHalf,DeligneWeilI,DeligneWeilII}. If a weight-$w$ Frobenius is $q^{w/2}$ times a positive-metric unitary, each eigenvalue has modulus $q^{w/2}$. The norm conclusion is immediate; constructing that form or proving purity is the deep step. Symplectic/alternating duality $F^tJF=q^wJ$ alone only pairs reciprocal eigenvalues and does not give their individual moduli. An order resolution or Schur-positive character expansion alone cannot replace this step in number-field cohomology.

The recent primary comparisons used here are specific:
\begin{itemize}
\item \cite[Theorem3.12, pp.15--17]{ConnesConsaniKnots} gives finite abelian covers with local orbit length $\log p$ and genuine arithmetic monodromy. It suggests what a descent-equipped order functor should match; it does not identify our $u$ or a face cycle with a prime without such a functor.
\item \cite[Propositions4.1--4.3, Theorem4.6]{Prolate} builds finite-prime local models with archimedean Fourier factors. Its Hilbertian isomorphisms are not an asserted uniformly controlled common isometry for our arithmetic approximants. \cite[Proposition3.1]{LocalMoments} gives determinacy of a particular local moment problem; positive finite Hankel matrices are not the missing global Weil inequality.
\item \cite[Theorem3.6]{ZetaTriples} gives discreteness/lower boundedness for the stated restricted Weil form. Theorem5.10, pp.23--24, additionally assumes a finite truncation, $\lambda>1$, a simple even bottom eigenfunction $\xi$ and $\delta_N\xi=1$. Its quotient metric is positive and its regularized determinant is $-i\lambda^{-iz}\widehat\xi(z)$. Passing to the infinite arithmetic comparison, including the ground-state and convergence conditions in its \S8, is not automatic.
\item \cite[Theorem6.1, p.13]{QuadraticForms} assumes a specified real distribution on an interval, the corresponding trigonometric-core form, essential self-adjointness, and an isolated simple lowest eigenvalue with reflection-even eigenfunction. Its Fourier transform then has only real zeros. Our degree operator's vacuum and length parity are not those data.
\item \cite[Theorem6.1]{DeningerPackets} gives prime-indexed periodic packets in its dynamical construction, not one proved orbit per prime with the full desired cohomological trace formula. \cite{DeningerReal} isolates an obstruction to a particular real-coefficient cohomology with full functorial Artin multiplicities; it is not a theorem excluding every Hilbert--P\'olya operator. The current survey \cite{ConnesSurvey} distinguishes remaining steps. The recent Picard-monoid construction \cite{ConnesConsaniJacobian} is not already a polarized abelian variety with a number-field purity theorem. \cite{PongeTian} uses RH as an input in the cited spectral identification, so cannot be used circularly as its proof.
\end{itemize}
These are targeted theorem/proof readings, not claims that every page of every cited work has been audited. The preserved review records specify actual page coverage and hashes, including inaccessible original PDFs where the modern primary proof was used instead.

\section{A concrete sequence of remaining constructions}
First choose a number-field model of a Belyi pair, its finite covers and conductor/ramification data. Construct actions or exact correspondences on the labelled order/module or analytic categories that recover the local prime Frobenius polynomials. The elliptic twists show why omitting this datum loses information. Test at more than one arithmetic model of the same complex dessin, not just its topological passport.

Second specify a finite projective or otherwise determinant-admissible complex and its \emph{independent} arithmetic endomorphism. Prove exact-sequence/determinant-line compatibility. Use semi-invariants to encode presentation degeneracies, and Schur/exterior characters to calculate local factors, retaining every central-action and nilpotence distinction above. General derived Schur replacements require their own construction; termwise Schur on the resolution fails the exhibited split sequence test.

Third prove the same local factors, gamma terms and conductor give a globally convergent or correctly regularized determinant. Trace-norm convergence must be established for the chosen analytic perturbations; strong convergence is insufficient in general. The inserted Abel family supplies this topology for an elementary scalar encoding, but has neither the independent arithmetic spectrum nor the positive metric needed for zero location.

Finally produce an independently proved positive form or self-adjoint determinant-preserving realization. One acceptance criterion is the Hurwitz target in the core; another is equality with the precise arithmetic quadratic form in the recent spectral theorems, with all their domain, ground-state, normalization and limit hypotheses verified. The local Cauchy kernel, $K$-group computation and graph spectral dimension alone meet none of these final comparison hypotheses. No remaining condition is silently asserted true here. Dirichlet, Hecke, Dedekind and Artin scopes must remain individually recorded.


\chapter{Graph pullbacks and completions with specified fibers}
\section{The complete Fock graph presentation}
For graph C* conventions the edge operator has initial projection at its range vertex and final projection under its source vertex. Put
\[
 S_h=L_{a_h}Q_{a_{\tau h}},\qquad T_h=L_{a_h}v_{[\tau h]}.
\]
The two right-hand projections are orthogonal subprojections of $L_{a_h}^*L_{a_h}=Q_{a_{\tau h}}+v_{[\tau h]}$. Thus $S_h^*S_h=Q_{a_{\tau h}}$, $T_h^*T_h=v_{[\tau h]}$, and $S_h^*T_h=0$. For distinct $h$ the prefix ranges are orthogonal, so every other cross-edge CK1 product is zero. Their final projections sum to $Q_{a_h}$. The vacua have no outgoing edges. The resulting graph has $|H|+|Q_0|$ vertices and $2|H|$ edges, with cycles only among its state vertices and exits to the paired sinks. Gauge covariance and nonzero vertices prove faithfulness as well as surjectivity. This explains precisely how a graph C* algebra is obtained without imposing the incompatible original-quiver CK1 relations.

The graph integer map has one column for each regular vertex $h$ and no sink columns:
\[
 d:\Z^H\longrightarrow\Z^H\oplus\Z^{Q_0},\qquad
 d(e_h)=e_h-e_{\tau h}-b_{[\tau h]}.
\]
The state part of $d(x)=0$ says $x_h=x_{\tau^{-1}h}$; the sink part then says the constant cycle coefficients are in $\ker B_\Gamma$. For the cokernel choose a base state on each cycle and successively use $q_h=q_{\tau h}+b_{[\tau h]}$ to eliminate all other states. Closing the cycle leaves $\sum_{h\in c}b_{[h]}=0$, while its base state remains free. Hence
\[
 \operatorname{coker}d\cong\Z^{\#\operatorname{cyc}\sigma}\oplus\operatorname{coker}B_\Gamma.
\]
This gives an independent integral check of the previous compact-extension computation. The full labelled graph, diagonal projections and edge families carry more information than those coarse groups.

\section{A genuine application of the supplied pullback theorem}
For the interval dessin let $E_2=G_\alpha$ have loop vertices $q_x,q_y$, common sink $b$, loops $s_x,s_y$ and exits $t_x,t_y$. Remove $b$ to obtain $F_2$, two disjoint loop graphs. Its quotient is exactly $C(\mathbb T)^2$. Let $E_1$ have the same three vertices and edges $e_{x,n}:q_x\to b$, $e_{y,n}:q_y\to b$, $n\ge0$, and no other edges. Remove $b$ to obtain the edgeless $F_1$, whose algebra is $\C^2$.

The inclusions $F_i\subset E_i$ are admissible in the precise sense of \cite[Definition3.1]{GraphPullbacks}: the removed sink set is hereditary/saturated; all edges ending at retained vertices are retained; no vertex emits infinitely many edges into the removed set and a positive finite number outside it. For $E_1$ the latter outside count is zero, which is allowed. Define a path functor identity on objects by
\[
 e_{x,n}\longmapsto s_x^nt_x,\qquad
 e_{y,n}\longmapsto s_y^nt_y.
\]
Images of distinct positive edges are prefix incomparable: at the first different length one has an exit and the other a loop. Source vertices and identities cause no exception. The image consists of all positive pointed paths in $E_2$, whose last edge is not a short loop; vertices are included separately. The first graph is acyclic and the removed sink has no short loop in the second. All hypotheses of Theorem3.4 are now verified, giving
\[
 C^*(E_1)\cong\mathcal T_A\mathop{\times}_{C(\mathbb T)^2}\C^2.
\]
The quotient of $C^*(E_1)$ by its sink ideal is $\C^2$, and its sink ideal is compact operators on the paths ending at $b$. The graph is acyclic, hence its algebra is AF. There are no regular vertices in $E_1$: the two sources are infinite emitters and $b$ is a sink. The graph $K$ matrix has no columns and three rows, yielding $K_0=\Z^3$ and $K_1=0$.

This is a positive instance of the new source inside our infinite dessin framework. It is a constants-symbol subalgebra of the Fock model, not the entire Fock algebra or the formal order. For a permitted cycle of length greater than one, its full closed path is itself pointed in this paper's convention, so an object-preserving acyclic path functor cannot cover all such paths. That simple obstruction prevents copying this example to every cycle by analogy. The sink-extension Theorem4.3 also requires a specified original pullback, injective gluing maps, annihilation of glued vertex projections and coverage of all paths ending at the gluing sinks. A two-to-one identification of sink occurrences within one branch does not satisfy these assumptions without another construction.

\section{Diagonal residue obstruction and the corrected pullback}
Suppose one tries to glue full matrix-circle blocks by their diagonal values. In $M_2$, $E_{12}E_{21}=E_{11}$, although both factors have zero diagonal. Thus the proposed boundary map is not a homomorphism. For $\ell>1$, more generally any *-homomorphism from the full matrix-circle algebra to a commutative C*-algebra has equal images of $E_{ii}=E_{ij}E_{ji}$ and $E_{jj}=E_{ji}E_{ij}$. These images are both orthogonal and equal, hence zero; summing the diagonal units kills the image of one and then the whole homomorphism. This works for nonunital *-maps as well: a nonzero image of one would give a nonzero unital corner map. The failure concerns full matrix-circle blocks, not the already-restricted fiber algebra $D_\ell$.

Evaluation $D_\ell\to\C^\ell$ is a genuine surjective *-homomorphism since at one both factors are diagonal. Its kernel is $C_0(\mathbb T\setminus\{1\},M_\ell)$. Pulling it back along the repeated-coordinate embedding $\Delta_\alpha:\C^{Q_0}\to\C^H$ gives exactly the fiber conditions of $\mathcal D_H$. The paired constant diagonal matrices give a unital *-section. Hence this is a split C*-extension, not a claim deduced from a merely commutative diagram of nonmultiplicative maps. Since $\mathbb T\setminus\{1\}\cong\R$, each kernel block is a matrix suspension of $\C$, with $K_0=0$, $K_1=\Z$. Split exactness now proves Theorem~\ref{thm:diagonalmodel}.

The pairing is retained in this specified presentation, not recovered from the unmarked isomorphism class. In particular arbitrary constant permutation matrices can relabel a branch's diagonal positions, so the spherical and elliptic theta examples have isomorphic unmarked $\mathcal D_H$ despite different ribbon orderings. Keep the labelled arrows, cycle cuts and distinguished fiber as additional data when using it as a dessin model. Also the central branch unitary $zI_\ell$ represents $\ell$ times the primitive corner generator in $K_1$, not a new independent generator per matrix entry.

\section{Generation, faithfulness and the exact failure of extension}
Let $\rho$ denote the polynomial representation in Theorem~\ref{thm:diagonalmodel}. Its vertex images are orthogonal projections summing to one. Every arrow image lies under the correct source and range vertex projections. A forbidden product vanishes because the intermediate state indices are $\tau h$ and $\alpha\tau h$, which are distinct. Thus $\rho$ is an actual algebra map, not a fitted characteristic polynomial.

For a full turn starting at $h$, $\rho(c_h)=w_cE_{hh}$ with $w_c(z)=g(z)^{|c|}z$. It is a normal element of $\mathcal D_H$. The function $w_c$ vanishes exactly at one; away from that point $w_c/|w_c|=z$. Therefore it is injective on all of $\mathbb T$, and a homeomorphism to its image. Nonunital continuous functional calculus at $\rho(c_h)$, using functions zero at zero, gives every $fE_{hh}$ with $f\in C(\mathbb T)$, $f(1)=0$. The zero condition makes the result vanish in all other matrix entries and branches.

To obtain an off-diagonal entry choose the unique shortest permitted path from state $h$ to $j$ within its cycle, of length $r>0$ and cut-crossing count $\kappa\in\{0,1\}$. Its image is $g^rz^\kappa E_{hj}$. Given $f(1)=0$, multiply it on the left by the already available diagonal element
\[
 \frac{f(z)z^{-\kappa}}{g(z)^r+\varepsilon}E_{hh},\qquad \varepsilon>0.
\]
The product is $f g^r/(g^r+\varepsilon)E_{hj}$ and converges uniformly to $fE_{hj}$. Indeed near one $|f|$ is small, and on the complement of that neighborhood $g^r$ has a positive minimum. These two bounds prove uniform convergence without assuming division by a function vanishing at the fiber. Thus the generated algebra contains the whole kernel ideal. Its quotient contains the entire paired diagonal algebra, so it is $\mathcal D_H$.

For injectivity let a finite linear combination of path images vanish. At one, all positive paths vanish, forcing the coefficients of the original vertex idempotents to vanish. Fix a state-matrix entry in a branch. Its remaining paths have lengths $r+m\ell$ and images
\[
 g^rz^\kappa\sum_m a_m(g^\ell z)^m E_{hj},
\]
with a possible initial zero coefficient if $r=0$. On $\mathbb T\setminus\{1\}$ the prefactor is nonzero, and $w_c$ has infinite range, so the polynomial has infinitely many distinct zeros. All coefficients are zero. Repeating by entries proves faithfulness.

No boundedness improvement repairs full formal extension. Pick any $z_0\ne1$, let $\lambda=w_c(z_0)^{-1}$, and consider $1-\lambda c_h$. Its inverse $\sum_{m\ge0}\lambda^m c_h^m$ exists in the arrow-adic order. The matrix image at $z_0$ has a zero eigenvalue on state $h$, so cannot be invertible. A unital algebra homomorphism would preserve the inverse identity, an immediate contradiction. This proves nonextension without any continuity assumption. It does not rely on the false claim that a norm greater than one alone forces a geometric-series denominator to be singular. Scaling $g$ to have norm at most one still does not allow arbitrary formal series.

The two valid operator completions are structurally different. In the bipartite case $K_0(\mathcal T_A)=\Z^{E+1}$ differs from $K_0(\mathcal D_H)=\Z^E$; in the nonbipartite case the former has the extra $\Z/2$ torsion. Neither is automatically the norm completion of the arrow-adic topology. No invariant constructed solely from these complex objects distinguishes the rational twists already exhibited; arithmetic descent remains a separate input.

\section{Invariant basis number and the Leavitt comparison}
If a unital $k$-algebra $R_1$ has a unital representation on a nonzero finite-dimensional vector space $V$, an isomorphism $R_1^m\cong R_1^n$ induces an isomorphism $V^m\cong V^n$ by tensoring. Comparing dimensions gives $m=n$. This proves invariant basis number without requiring the representation to be faithful. Every original dessin algebra and completed order has a vertex-simple representation, and each permitted-state Laurent algebra has evaluation onto a nonzero matrix block. The Leavitt algebra of $G_\alpha$ has the same finite-dimensional representation after killing the vacuum sinks and evaluating one boundary cycle. Thus all of these examples have invariant basis number. Leavitt's free-rank phenomena are important elsewhere but cannot be assumed in this class simply from the name LPA.

For the specified finite graph $G_\alpha$, the usual complex Leavitt path algebra embeds in $C^*(G_\alpha)$ with norm-dense image, using its canonical generators and involution \cite[Theorem7.3]{TomfordeEmbedding}. Those generators are the split $S_h,T_h$, not the unsplit original gentle arrows equipped with ordinary CK1. This is the exact algebraic/analytic twin relationship applicable here.


\chapter{String finiteness and the smooth quotient surface}
For a finite quiver, finitely controlled implies finitely generated: choose one finitely generated submodule containing each vertex space and sum the finitely many choices. Conversely the module itself works. Stabilization at each of finitely many vertices likewise gives ordinary artinianity. Crawley--Boevey's polynomial theorem therefore applies directly to our finitely generated and artinian polynomial modules, with its primitive injective bands in the latter case. It does not assert a classification of every possible algebraic module.

The formal branch $M=k[[x]]$, $y=0$, gives a precise counterexample to that stronger assertion. An $A$-endomorphism $F$ preserves each $x^nM$. For the polynomial truncation $f_n$ of $f$, one has $F(f)-f_nF(1)\in x^nM$. Separatedness and all $n$ imply $F(f)=fF(1)$. Hence $\operatorname{End}_A(M)=k[[x]]$ is local, so $M$ is indecomposable. If an ordinary word module has $y=0$, its word cannot contain $y$ or $y^{-1}$, since such a letter would create a nonzero basis map. A reduced word using only $x^{\pm1}$ cannot change direction. Its nonperiodic choices are finite strings, the free ray $k[x]$, and the opposite locally nilpotent ray. The first and last fail injectivity of $x$, and the free ray fails invertibility of $1-x$. A two-sided pure-$x$ word is periodic, and in every associated band $x$ acts invertibly. On $M$, $x$ is injective but not onto. None of these possibilities equals $M$. This includes arbitrary nonzero indecomposable Laurent coefficient modules, not just finite bands.

For a general nonzero completed lattice, finite generation over $A$ would imply finite generation over $k[u]$. Every nonzero polynomial in $u$ is a power of $u$ times a unit of $k[[u]]$, so it acts injectively. The lattice would be a nonzero finite free $k[u]$-module. Multiplication by $1-u$ on such a free module has a nonzero cokernel, but the same map is invertible over $k[[u]]$. This proves the contradiction in the core. The strictly descending $u$-adic chain separately excludes artinianity. These are actual failed hypotheses, not merely unverified applicability.

Lekili--Polishchuk's precise convention takes $D(A)$ to mean perfect dg modules. Their locally gentle definition allows infinite dimension, and their smoothness criterion says no forbidden cycle; properness says no permitted cycle. Both fail for every nonempty closed dessin here. Their Theorem3.2.2 therefore does not apply to the full polynomial dessin algebra, even though it applies to suitable infinite-dimensional quotients. The OPS finite-dimensional geometric-intersection theorem has its own grading, local-system, boundary and exceptional coincident-curve conditions. An unrestricted statement that every gentle/string derived object is an ungraded geodesic crossing is not justified by either source.

For the interval quotient $B=A/(y)=k[x]$, global dimension one follows from its PID structure: every submodule of a finite free module is free, giving finite projective resolutions of all finitely generated modules. Thus $\operatorname{Perf}(B)=D^b(\operatorname{mod}_{\rm fg}B)$ in grading zero. Its permitted loop gives one unmarked boundary component and its length-one forbidden thread $x$ gives one stop on the other boundary, the annulus of the smooth partially wrapped model. This construction uses the infinite polynomial quotient, not the completed $k[[x]]$ category.

Finally the quotient does not give a derived fully faithful embedding by ordinary inflation. Resolve $B$ over $A$ by the alternating maps $y,x,y,x,\ldots$. Hom into $B$ gives $0,x,0,x,\ldots$, so $\operatorname{Ext}_A^{2r}(B,B)=B/xB=k$ for every $r\ge1$, whereas all positive self-Ext over $B$ vanish. Consequently its inflation is not derived fully faithful, and the perfect $B$-module $B$ becomes an infinite-projective-dimension $A$-module. A surface functor transferring these categories must account for that difference explicitly.



\chapter{Arithmetic character sums and the cellular determinant}
\section{The uniform calculation and its exceptional cases}
In the finite-field calculation write $k=\mathbb F_q$; the final order-insertion paragraphs specify a separate characteristic-zero coefficient field. The curve is nonsingular because its discriminant is $-432d^2\ne0$, and has one point at infinity. Thus counting $1+\eta(x^3-d)$ in each affine fiber proves the trace formula in Theorem~\ref{thm:jacobitheta} directly. The multiplicative group of $k$ is cyclic. When $3\mid q-1$, a nonzero cube has three roots and a noncube none, detected by $1+\rho(t)+\bar\rho(t)$. At zero this expression is one, using the \emph{constant} one. This convention differs from mathlib's zero-extended trivial multiplicative character, whose value there is zero. The proof and formal components never replace the constant by that character.

The substitution is valid including $u=0,1$ because the characters are zero extended:
\[
 \sum_t\rho(t)\eta(t-d)
 =\sum_u\rho(du)\eta(du-d)
 =\rho(d)\eta(-d)\sum_u\rho(u)\eta(1-u).
\]
Multiplying by the initial minus sign gives exactly $\pi_d+\bar\pi_d$. Changing $\rho$ to $\rho^{-1}$ conjugates both $j$ and $\lambda_d$, so does not change the unordered roots or the integer trace. Characteristics two and three and $d=0$ are excluded throughout. For a residue characteristic $p\equiv2\pmod3$, odd extension degrees have trace zero by cube bijectivity, whereas even degrees lie in the split case; no claim of zero trace for all extensions of that characteristic is made.

For completeness, choose a nonzero $\mathbb F_p$-linear functional on $k$ and compose it with $t\mapsto e^{2\pi it/p}$ to obtain $\psi$. For a nontrivial multiplicative character $\chi$,
\[
 \begin{aligned}
 |G(\chi)|^2
 &=\sum_{u\ne0}\chi(u)\sum_{y\ne0}\psi((u-1)y)\\
 &=(q-1)-\sum_{u\ne0,1}\chi(u)=q.
 \end{aligned}
\]
Both orthogonality identities follow by translating a finite character sum by an element on which its character is not one. The $z=0$ terms in $G(\chi)G(\theta)$ sum to $\theta(-1)\sum_x(\chi\theta)(x)$, zero if the product is nontrivial. For $z\ne0$ set $x=zu$, $y=z(1-u)$ to obtain $J(\chi,\theta)G(\chi\theta)$. Coprime orders three and two ensure our product character is nontrivial. Taking norms proves $|J|^2=q$, rather than presuming a desired positivity. This is the classical calculation reviewed directly in \cite{WeilJacobi}; no claim of novelty is made.

At seven the five nonzero summands are $-\omega^2,-\omega,\omega,-\omega^2,1$, so $J=3+2\omega$. Multiplication by $\lambda_2=\omega^2$ gives $-1-3\omega$. Their matrices are
\[
 \begin{pmatrix}3&-2\\2&1\end{pmatrix},\qquad
 \begin{pmatrix}-1&3\\-3&2\end{pmatrix}.
\]
Both have determinant seven and the same positive Gram form; their traces are four and one. The arithmetic distinction between the rational models now comes from explicit finite-field characters, while their complex dessin is unchanged. Exact prime-field diagnostics test every nonzero twist for the listed nineteen primes; the theorem itself is uniform in prime powers and rests on the proof, not that finite list.

\section{Deck descent and the actual cellular basis}
The CM cellular construction in this and the next subsection assumes $q\equiv1\pmod3$.
The cover $\beta=x^3/d$ has deck transformations $(x,y)\mapsto(\zeta_3x,y)$ and $(x,y)\mapsto(x,-y)$, generating $C_3\times C_2=C_6$. Label their actions on darts by $\sigma=(012)(345)$ and $\alpha=(03)(14)(25)$; they commute. Their product is $(0\,4\,2\,3\,1\,5)$. They preserve the edge pairing and both successor permutations. Consequently they define algebra automorphisms, preserve $J$ and its completion and fix the common $u$. The projective-resolution maps are compatible since $\theta_g(a_hx)=a_{gh}\theta_g(x)$. Twisting module actions gives an exact functor and the corresponding relabelled resolution, without an assumption of arithmetic Frobenius.

Over an algebraic closure choose $a^3=d$, $b^2=d$ and $\varphi(x,y)=(ax,by)$ from $E_1$ to $E_d$. For the $q$-power Frobenius morphisms,
\[
 \varphi^{-1}\circ F_d\circ\varphi=h_d\circ F_1,\qquad
 h_d(x,y)=(a^{q-1}x,b^{q-1}y).
\]
When $q\equiv1\pmod6$, the first coordinate is the cubic residue $d^{(q-1)/3}$ and the second $\eta(d)$. Choose the embedding of cube roots defining $\rho$ consistently. On the deck eigenspace labelled by $dx/y$, the \emph{deck} character is their ratio, $\rho(d)/\eta(d)=\lambda_d$. This labels the twist, not a nonzero Frobenius action on differentials in characteristic $p$: that pullback is zero. An inverse convention for arithmetic Galois Frobenius changes the action convention; our displayed formula uses the geometric curve morphism explicitly. The character-sum polynomial fixes its trace without an unstated convention.

Orient the three geometric edges as in the core. Each occurs twice with opposite orientation in the single face, so $\partial_2=0$; $\partial_1$ has columns $(-1,1)^t$. The homology basis $b_0=e_0-e_1$, $b_1=e_1-e_2$ has Gram matrix $G$ inherited from the orthonormal cellular edges. Rotation sends $b_0\mapsto b_1$, $b_1\mapsto-b_0-b_1$, exactly $C$. The involution swaps vertices and reverses each edge, hence acts as $-I$ on this plane. Thus multiplication by $A+B\omega$ is literally $AI+BC$, not merely an abstract matrix with the correct trace. Expanding its determinant gives the norm $A^2-AB+B^2$, and multiplication of this norm proves $M^tGM=qG$. The displayed positive quadratic decomposition proves the metric independently of unknown zeros.

\section{An explicit chain lift and determinant cancellation}
Let $P$ be the three-cycle permutation on $C_1$. Write $\pi_d=A+B\omega$ and $\gamma=A+B$. Define
\[
 F_1=AI+BP,\qquad
 F_0=\frac12\begin{pmatrix}1+\gamma&1-\gamma\\1-\gamma&1+\gamma\end{pmatrix},\qquad
 F_2=q.
\]
These are rational cellular maps. The sum-zero plane of $F_1$ has restriction $M_d$, while the edge-sum line has eigenvalue $\gamma$. On $C_0$ the vertex-sum line has eigenvalue one and the boundary line eigenvalue $\gamma$. It follows directly that $F_0\partial_1=\partial_1F_1$; the other chain condition is zero. The induced maps on $H_0,H_1,H_2$ have characteristic factors $1-T$, $P_d(T)$, $1-qT$. In particular
\[
 \det(I-TF_0)=(1-T)(1-\gamma T),\qquad
 \det(I-TF_1)=(1-\gamma T)P_d(T).
\]
The contractible boundary factor cancels, giving
\[
 \frac{\det(I-TF_1)}{\det(I-TF_0)\det(I-TF_2)}
 =\frac{P_d(T)}{(1-T)(1-qT)}.
\]
This equality holds as a formal rational identity and by cancellation at apparent intermediate singularities. The alternating determinant with exponent $(-1)^i$ is its inverse; signs are not changed silently. The standard elliptic trace and determinant theorem identifies the right side with the finite-field curve zeta function \cite{SGAFourHalf,DeligneWeilI}. The elementary proof above supplies its actual trace and norm here; it does not independently reprove the general cohomological trace formalism.

This lift also has the required descent character on $H_1$. Write the base Jacobi value as $j=A_1+B_1\omega$. Then $h_d$ acts there by $\eta(d)C^r$ when $\rho(d)=\omega^r$, so twisting $A_1I+B_1C$ multiplies its eigenvalue $j$ by $\lambda_d$, giving $M_d=A_dI+B_dC$. Reduction modulo $C^2+C+I=0$ gives the coefficients of $\pi_d$. If the corresponding full three-edge polynomials differ by $c(I+P+P^2)$, the coefficients are $c=0,\eta(d)B_1,\eta(d)A_1$ for $r=0,1,2$ respectively; they have the same homology action and differ only on the contractible boundary pair. An explicit homotopy sends $v_--v_+$ to $c(e_0+e_1+e_2)$ and the vertex sum to zero. Its $\partial H+H\partial$ is the edge difference and the boundary scalar difference $3c$. Thus this normalization is compatible up to rational chain homotopy. It is not asserted to arise from a continuous Frobenius map of the complex Riemann surface, or to give a canonical integral Tate-lattice comparison.

\section{A decisive restriction at inert primes}
For $q\equiv2\pmod3$, Frobenius sends $\zeta_3$ to $\zeta_3^{-1}$. Thus its cohomological relation with the cubic deck action requires $MC=C^2M$ (the pullback convention yields the equivalent reversed relation). On a real two-dimensional carrier with our fixed $C$, comparing entries forces
\[
 M=\begin{pmatrix}a&b\\a+b&-a\end{pmatrix},\qquad
 \det M=-a^2-ab-b^2
 =-\tfrac12\bigl(a^2+b^2+(a+b)^2\bigr)\le0.
\]
The genuine elliptic Frobenius determinant is $q>0$, so such a real carrier is impossible. This is an elementary instance of the real-coefficient cohomology issue discussed by Deninger \cite{DeningerReal}. It does not rule out a complex realization, a larger real space, a different coefficient theory or an infinite-dimensional height operator. In particular it is not a refutation of an unrestricted Hilbert--P\'olya approach.

There is a local complex replacement, with its limitations explicit. The conjugation matrix
\[
 K=\begin{pmatrix}1&-1\\0&-1\end{pmatrix}
 \quad\text{satisfies}\quad K^2=I,\ KC=C^2K,\ K^tGK=G.
\]
Set $F=i\sqrt q\,K$ on the complexification. Then $\tr F=0$, $\det F=q$, $F^2=-qI$ and $F^*GF=qG$, giving the correct trace-zero local factor. This is a two-dimensional complex, hence four-dimensional real, local model. This is a one-$q$ model; independently choosing $i\sqrt q K$ for each extension field does not give Frobenius-power compatibility. Such a family would need to start from one $F_p$ and prove the comparison for its powers. The inserted $i\sqrt q$ is not a canonical rational or integral Tate comparison or a construction of a global positive determinant. The fixed real two-dimensional model has genuinely failed; its surviving split-prime construction and this complex alternative must not be conflated.

\section{The remaining order and global spectral obligations}
The deck automorphism has finite order, so over characteristic zero its eigenvalues are roots of unity. Its permutation action alone cannot supply a weight-one norm $\sqrt q>1$. The independently computed Jacobi coefficient is essential. The above construction closes the specified local cellular benchmark, not a reconstruction of arithmetic from bare complex dessin data.

For these order-insertion tests take $\Lambda=\Lambda_{\mathbb{Q}}$ over $R=\mathbb{Q}[[u]]$, with rational residue simples and $V_d=\mathbb{Q}^2$. This is distinct from reducing the matrix in characteristic $p$, where its determinant $p$ becomes zero. There are two concrete failures of a naive order lift. First, a residue simple $S_e$ is annihilated by all arrows; $S_e\otimes V_d$ with independent endomorphism $I\otimes M_d$ has the desired underlying finite-$\mathbb{Q}$ determinant. But its projective dimension is infinite by the nonterminating forbidden resolution, and $\sigma$ permutes all three vertex labels, so choosing $S_e$ does not preserve full deck symmetry. This is not a perfect-order determinant realization.

Second, $\Lambda\otimes V_d$ is free of rank two over $\Lambda$, but the underlying $R$-rank of $\Lambda$ is eighteen. In an $R$-basis, the \emph{$\Lambda$-linear} $I\otimes M_d$ is eighteen copies of $M_d$, hence its determinant is $P_d^{18}$. On $e\Lambda$ the rank is six: each of its two hereditary branches contributes a row of length three, with finite conductor gluing changing no rank. This gives $P_d^6$. Restoring all three residue vertices similarly gives $P_d^3$ for this untwisted tensor action. A genuine deck-semilinear map permutes factors as well and need not have these powers; it requires its own comparison. No justified determinant-line root or multiplicity-removing functor is supplied here.

Finally fix a nonzero rational $d$ when forming its global Euler product from good reductions. Local purity cannot be multiplied into a global zero-location proof. Since each reciprocal root has modulus $\sqrt p$, the product of good elliptic factors is absolutely convergent on $\Rea s>3/2$. The local zeros at $\Rea s=1/2$ are outside that initial product argument. The finite-field curve zeta's center is $1/2$; the global elliptic $L$-function's unshifted center is one, or $1/2$ only after a half-unit shift. Neither is an identification with the primitive Dirichlet function targeted earlier. Ramified factors, archimedean normalization, adelic descent and the global determinant-preserving positive spectral limit remain separate obligations.



\chapter{The theta order interface beyond tensor insertion}
We prove the multiplicity theorem and construct its actual functorial escape, retaining the completed infinite order as the primary algebra. The finite algebra used below is its specified special fiber, not a replacement primary example.

\section{Normal forms and the special fiber}
Index the three quiver vertices by $i\in\mathbb Z/3$. Let $a_i^\pm:i\to i-1$ be the two permitted branches. Only paths of constant sign survive. Write $p_{i,r}^\pm$ for their length-$r$ paths and $c_i^\pm=p_{i,3}^\pm$. The split order model gives
\[
 ue_i=c_i^++c_i^-,\qquad
 p_{i,r+3}^\pm=u p_{i,r}^\pm\quad(r\ge1).
\]
The second identity follows because multiplying by the opposite full cycle switches branch and vanishes. An $R$-basis of $e_i\Lambda$ is
\[
 e_i,\ a_i^+,a_i^-,\ p_{i,2}^+,p_{i,2}^-,\ c_i^+.
\]
These span by the two identities and completion. They are independent in the two hereditary branch orders: off-diagonal entries distinguish the four shorter paths; the remaining combination $xe_i+yc_i^+$ has paired diagonal entries $(x+uy,x)$, forcing $x=y=0$. This is an actual power-series basis, not just a generic rank calculation.

In $B=\Lambda/u\Lambda$, the same eighteen basis vectors remain. A product of nonempty paths is zero if branches switch or the total length exceeds three; the length-three minus path is the negative of the plus path. A length-four path is $u$ times a length-one path, hence zero. The radical $\mathfrak r$ is the positive-length ideal, with $\mathfrak r^4=0$ and $B/\mathfrak r=\mathbb Q^3$. Each $e_iB$ has radical-layer dimensions $(1,2,2,1)$ and has two composition factors at each vertex, so the Cartan matrix is
\[
 C_B=\begin{pmatrix}2&2&2\\2&2&2\\2&2&2\end{pmatrix}.
\]
The Cartan matrix alone does not establish the endomorphism formula; the path filtration supplies that step.

\section{All projectives and all linear endomorphisms}
Let $J$ be the completed arrow ideal. Every $1-x$, $x\in J$, has its convergent formal inverse; the same applies to $1-yx$ for arbitrary $y\in\Lambda$. Thus $J$ is contained in the Jacobson radical, and $\Lambda/J=\mathbb Q^3$ implies equality. If $P$ is finitely generated projective, choose a basis of each multiplicity space in $P/PJ$ and lift its vectors into $Pe_i$. This gives
\[
 Q=\bigoplus_i(e_i\Lambda)^{n_i}\longrightarrow P
\]
inducing an isomorphism on tops. Nakayama makes this surjective. Its split kernel $N$ is finitely generated and satisfies $N/NJ=0$, hence vanishes by Nakayama. This proves the decomposition for every such projective without an unproved freeness assumption over $\Lambda$.

An entry $e_i\Lambda\to e_j\Lambda$ is left multiplication by an element of $e_j\Lambda e_i$. Its degree-zero part is possible only when $i=j$ and is a scalar multiple of $e_i$. For $f$ on $P$, write its top blocks $A_i$ and set $A=\operatorname{diag}_iA_i$. The remainder has positive path length. On $P/uP$ it strictly raises the path filtration. At each of the two length-one and two length-two slots, and the length-zero and length-three slots, the degree-preserving map is exactly $A$. In an increasing-degree basis the matrix is lower block triangular. For the formal upper-block convention order the six slots instead as
\[
 3+,\ 2+,\ 2-,\ 1+,\ 1-,\ 0.
\]
There are six diagonal copies of $A$, allowing arbitrary positive-length off-diagonal entries. Hence
\[
 \det_{\mathbb Q}(1-T\overline f)=\det_{\mathbb Q}(1-TA)^6.
\]
Specializing the finite $R$-matrix determinant coefficientwise gives Theorem~\ref{thm:theta-multiplicity}. No simultaneous diagonalizability, tensor form of $f$, arithmetic purity or deck equivariance is required.

For a bounded complex $P^\bullet$ with a linear chain map $f^\bullet$, fix either parity convention and write
\[
 D_R(P^\bullet,f^\bullet;T)=\prod_j\det_R(1-Tf^j)^{(-1)^j}.
\]
The inverse convention gives the reciprocal and the same obstruction. Each factor has constant coefficient one, so its reduction is nonzero. Specialization of this displayed rational expression is defined and equals
\[
 \left[\prod_{j,i}\det_{\mathbb Q}(1-TA_i^j)^{(-1)^j}\right]^6.
\]
Cancellations cannot change the divisibility of divisor multiplicities. In particular adding contractible projective pairs cannot repair it. Here a perfect endomorphism means a bounded-projective chain representative; no statement about an arbitrary infinite resolution or nonperfect lattice is inferred.

Generically put $N=\sum_i n_i$. As a right module over $M_3(K)\times M_3(K)$, $P_K$ is the sum of $N$ copies of each branch's three-dimensional row simple. Any commuting map has matrices $F_+,F_-\in M_N(K)$ on their multiplicity spaces, so
\[
 \det_K(1-Tf_K)=\bigl(\det(1-TF_+)\det(1-TF_-)\bigr)^3.
\]
For example, left multiplication by $c_i^+$ on $e_i\Lambda$ has ordinary characteristic determinant $(1-uT)^3$, so a sixth power is forced only at the special fiber. This also shows why generic reduced determinants and ordinary determinants must not be conflated. We do not need an integral choice of either generic matrix to derive the cube obstruction.

For the local numerator, suppose a sixth power $A(T)^6/B(T)^6$ equals a nonzero quadratic $P(T)$ with $B\ne0$. Polynomial degrees would give $6\deg A=2+6\deg B$, impossible. This includes repeated roots. For the complete curve zeta, $q>1$ and $|a_q|\le2\sqrt q$ imply $P(1)>0$ and $P(1/q)=P(1)/q>0$. Its two denominator roots are distinct uncancelled simple poles, again impossible for a sixth power. These arguments do not assume global RH or Artin holomorphy.

\section{An actual integral determinant functor of multiplicity one}
Project to one hereditary cycle order $B_+\subset M_3(R)$. Its standard column $W_+=R^3$ is a left $\Lambda$-module with $e_iW_+=R\epsilon_i$. More explicitly $a_i^+$ acts by $u^{w_i}E_{i,i-1}$, with $(w_0,w_1,w_2)=(1,0,0)$ for one cut, while every $a_i^-$ acts by zero. These matrices respect travel-order multiplication as written: $E_{ij}E_{jk}=E_{ik}$. Each plus full cycle is $uE_{ii}$, so the common central $u$ acts by $uI$. Nonempty paths tend to zero $u$-adically, giving the action of the \emph{full completed order}, not just its polynomial subalgebra.

This $(\Lambda,R)$-bimodule defines $-\otimes_\Lambda^{\mathbf L}W_+$. On a bounded complex of finite projectives it is computed termwise, and its terms are finite free over $R$. An acyclic bounded projective complex is split contractible, and tensor preserves that contraction. Thus it descends to a genuine triangulated functor $\Phi_+:\operatorname{Perf}(\Lambda)\to\operatorname{Perf}(R)$. No assertion of exactness on all un-derived modules or flatness of $W_+$ is needed. In fact $W_+$ has infinite projective dimension as a left $\Lambda$-module: $\Lambda$ is noetherian, being finite over the central noetherian ring $R$, and $W_+$ is finitely generated. A finite projective resolution could thus be chosen with finite terms. Every left projective term has $R$-rank divisible by six by the same top argument, whereas $W_+$ has rank three, contradicting the alternating rank identity. This does not prevent tensoring bounded right-projective complexes.

The isomorphism $e_i\Lambda\otimes_\Lambda W_+\to e_iW_+$ sends $e_i\lambda\otimes w\mapsto e_i\lambda w$. Its inverse is $e_iw\mapsto e_i\otimes e_iw$, proving rank one directly. Consequently $e_i\Lambda\otimes_{\mathbb Q}V$, with action $1\otimes M$, maps to $R\otimes_{\mathbb Q}V$ with action $M$. Its characteristic determinant is $\det(1-TM)$ \emph{once}. On the regular input, the three vertex lines survive and the factor appears three times. A full symmetric vertex orbit therefore is not silently replaced by a chosen vertex.

Different cycle cuts give different integral column lattices. Their generic modules are isomorphic, and branch endomorphism characteristic determinants agree by generic conjugation; equality over $K$ implies equality over $R$. Nevertheless an equivariant lattice is stronger data. Branch selection also changes the functor: for the unit $f=e_i+c_i^+$ on $e_i\Lambda$, the plus determinant is $1+u$ and its characteristic factor is $1-(1+u)T$. Under $\alpha$, $c_i^+$ becomes $c_i^-$, whose plus action is zero; the determinant becomes one and the characteristic factor $1-T$. Since these scalar automorphisms differ, no natural isomorphism of determinant lines can make this chosen branch functor $\alpha$-invariant for all endomorphisms. Pairing both branch functors makes this branch-swap data available but doubles the local factor on the basic-projective tensor input.

\section{The exact rotation obstruction and its allowed escape}
Consider any left $\Lambda$-lattice supported on the plus branch, free of rank one at every $e_i$, and with central $u$ acting as multiplication by $u$. Write its plus arrow scalars as $a_0,a_1,a_2\in R$. Their product is $u$, so they are nonzero and their nonnegative valuations sum to one. A strict $R$-linear lift of $\sigma$ that intertwines the order rotation permutes the three vertex lines by units. Its arrow-intertwining identities make the three arrow valuations equal. Thus $3m=1$ for an integer $m\ge0$, impossible. The argument already excludes an invertible intertwiner; imposing that its cube be one cannot help.

There is an elementary coefficient proof of the same obstruction. Intertwining by units makes vanishing of the three arrow constant coefficients equivalent. Their product has constant coefficient zero, so all three vanish. Factoring $u$ from two of them would make $a_0a_1a_2$ divisible by $u^2$, contradicting that it equals $u$. This is the literal power-series component checked in Lean.

After the explicit ramified change $R'=\mathbb Q[[v]]$, $u=v^3$, \emph{choose a new column lattice} with all three plus arrow maps $vE_{i,i-1}$. Cyclic permutation of the basis is an order-three intertwiner fixing $v$. The full-cycle relation is now $v^3=u$ as required. Merely base-changing the old cut lattice gives arrow valuations $(3,0,0)$ and does not produce that linearization. This construction changes the base and lattice and does not establish an unramified descent comparison. It also still needs both branches to address $\alpha$.

The rotation obstruction is about actual central-fixed rank-one branch lattices, not every possible determinant-level natural transformation. The branch characteristic determinants may be rotation-invariant through their generic reduced interpretation even without a lattice lift. Neither this theorem nor the ordinary-determinant obstruction excludes semilinear correspondences, other determinant categories or infinite analytic models.

The next arithmetic task is therefore narrower and more concrete: choose a branch functor, a ramified lattice or a different equivariant replacement only after proving the required descent comparison, and construct its prime action from arithmetic data rather than feed in the already known $M_d$. Bad primes, integral/Tate comparison, local-to-global conductor and gamma factors, and a positive spectral limit for the same completed function remain unresolved.



\chapter{Explicit tame descent and the full-deck coefficient test}
The preceding ramified column is a precisely specified object of the original infinite theta-order framework. We compute its descent instead of treating a semilinear group action as automatically effective integral descent. All constructions below are over the new \emph{mathematical} coefficient field $k=\mathbb Q(\omega)$; they do not replace the original rational model or its source history.

\section{The action on genuine formal series}
Write $R=k[[u]]$, $S=k[[v]]$ with $u=v^3$. Every $f(v)\in S$ has a unique expansion $f_0(u)+vf_1(u)+v^2f_2(u)$. Hence $S$ is finite free of rank three over $R$. Let $g(f)(v)=f(\omega v)$. On $W=S^3$ with basis $\epsilon_0,\epsilon_1,\epsilon_2$, put
\[
 \Theta(f_i\epsilon_i)=\omega^{-i}f_i(\omega v)\epsilon_i,\qquad
 C\epsilon_i=\epsilon_{i+1}\quad(i\bmod3).
\]
The plus arrow $a_i^+$ acts by $vE_{i,i-1}$, the minus arrow by zero. The cyclic product is $uE_{ii}$ and arbitrarily long permitted paths tend to zero $v$-adically, so this is a module over the full completed order. For $i=0$, the coefficient in the target under $\Theta a_i^+$ is $\omega v$ and in the source under $a_i^+\Theta$ it is $v\omega^{-2}$; these agree. The same calculation for $i=1,2$ proves compatibility with every arrow and therefore the completed algebra. Also $\Theta^3=1$.

For all coefficients, not just a truncated diagnostic,
\[
 [v^n](\Theta w)_i=\omega^{n-i}[v^n]w_i.
\]
Thus invariance is equivalent to support in $n\equiv i\pmod3$, and
\[
 L=W^\Theta=\bigoplus_{i=0}^2 Rb_i,\qquad b_i=v^i\epsilon_i.
\]
The arrows on this basis have scalar valuations $(1,0,0)$, since $a_0^+b_2=ub_0$, $a_1^+b_0=b_1$ and $a_2^+b_1=b_2$. Tame invariants return the old cut column rather than a new strict unramified rotation carrier.

The extension $R\subset S$ is faithfully flat. Its group quotient has a fixed point at $v=0$, however, and is not a $\mu_3$-torsor there: $S\otimes_R S$ is not the product of the three group graphs at that point. Semilinear $\mu_3$ data are therefore not by themselves the full effective fpqc descent data for a vector bundle on $\operatorname{Spec}R$. The adjunction comparison is explicitly
\[
 S\otimes_R L\longrightarrow W,\qquad
 \operatorname{diag}(1,v,v^2),
\]
with determinant $v^3=u$, cokernel $S/(v)\oplus S/(v^2)$ and closed-fiber rank one. This proves its failure to be an isomorphism directly. The quotient stack $[\operatorname{Spec}S/\mu_3]$ retains the stabilizer at zero; its coarse space is $\operatorname{Spec}R$. This local root construction is described in \cite{RootStackGuide}. Alper's vector-bundle criterion \cite[Theorem10.3, pp.2386--2387]{AlperDescent} distinguishes precisely this nontrivial closed-fiber action from effective coarse descent. We use these sources for interpretation; all displayed module comparisons have already been calculated explicitly.

\section{A projective symmetry survives on the endomorphism order}
Directly on the three basis vectors,
\[
 \Theta C=\omega^{-1}C\Theta=\omega^2C\Theta.
\]
Thus the rotation does not preserve $L$. The corrected $vC$ does preserve it, since the factor $\omega$ from $g(v)$ cancels the commutator. In the invariant basis it has matrix $A$ from the core, with $A^3=uI$ and $\det A=u$. This is generically invertible and gives an order-three projective rotation after passing to the fraction field; it is not an invertible integral lattice intertwiner.

There is nevertheless a true algebra-level escape. Let $F=(f_{ij})$ be an arbitrary $S$-linear endomorphism of the \emph{carrier} $W$, not assumed $\Lambda$-linear. Its tame conjugate has entries $\omega^{j-i}g(f_{ij})$. Invariance therefore means
\[
 f_{ij}\in v^{[i-j]_3}R,
\]
where $[r]_3\in\{0,1,2\}$ is the least residue. After replacing $\epsilon_i$ by $b_i=v^i\epsilon_i$, the entry becomes
\[
 v^{[i-j]_3+j-i}R=
 \begin{cases}R&i\ge j,\\uR&i<j.\end{cases}
\]
Consequently $\operatorname{End}_S(W)^\Theta$ is exactly the hereditary three-cycle order $B_+$ in the previously fixed cut convention. Multiplication is preserved by conjugation in the generic matrix algebra; no entrywise projection is being misused as a ring map.

The scalar commutator cancels under conjugation. Hence $\operatorname{Ad}(C)$ commutes with tame conjugation, descends to $B_+$, and is represented there by $\operatorname{Ad}(A)$ because $v$ is central. Since $A^3=uI$, this algebra automorphism has order three and preserves the hereditary order despite $A^{-1}$ not lying in its unit group. This explains the surviving rotation of the branch order; it does not construct a rotated integral column or a Frobenius endomorphism. The other branch is identical and the paired-diagonal residue gluing still defines the original $\Lambda_R$ inside $B_+\times B_-$. Its label automorphisms remain the already constructed dessin symmetries.

\section{Determinants, Schur weights and the boundary correction}
On $\det_S W=\bigwedge_S^3W$, both $\det D=1$ and $\det C=1$. Thus this line has the scalar action $g$ and genuinely descends to the trivial $R$-line with trivial rotation. This is a real positive determinant-level result. On the other hand,
\[
 b_0\wedge b_1\wedge b_2=v^3(\epsilon_0\wedge\epsilon_1\wedge\epsilon_2),
\]
so the natural comparison
\[
 \det_R(W^\Theta)\longrightarrow(\det_S W)^\Theta
\]
is multiplication by $u$, with cokernel $R/(u)$. Both lines are abstractly free of rank one, but their given comparison is not an isomorphism. It records the divisor of the ramified lattice change and cannot be discarded when claiming a determinant-preserving integral comparison. For the \emph{same} operator $vC$, the scalar on both determinant lines is $u$; naturality is respected. Comparing $\det C=1$ with $\det A=u$ without the extra factor $v$ would instead compare different operators.

For an induced nonzero homogeneous tensor/Schur carrier of degree $d$, the same functorial calculation gives commutator $\omega^{-d}$. The two induced invertible operators commute exactly when $3\mid d$. This is a group-action statement, not sufficient for coarse vector-bundle descent: $\operatorname{Sym}^3(S^3)$ has closed-fiber inertia multiplicities $(4,3,3)$ for the three characters, so its inertia is still nontrivial even though the commutator vanishes. A mixed tensor $W^{\otimes a}\otimes(W^*)^{\otimes b}$ has net exponent $a-b$; $\operatorname{End}_S(W)$ is the concrete net-zero escape computed above.

These carrier functors are not automatically functors on unital $\Lambda$-modules. Applying global $\bigwedge^3$ to the rank-one vertex idempotents kills all three, although their sum is the identity before applying this nonadditive functor. Vertexwise $\bigwedge^3$ of the rank-one column is zero. Vertexwise $\operatorname{Sym}^d$ for $d>0$ does retain the monomial zero relations but changes the full-cycle scalar from $u$ to $u^d$. It leaves the fixed central-action problem unless that change of base is included. None of the Schur cancellations supplies the missing two-dimensional arithmetic carrier.

\section{Why the local Jacobi plane does not descend by this repair}
On constant vectors $\Theta$ acts by $D$. It sends the sum-zero cycle $\epsilon_0-\epsilon_1$ to $\epsilon_0-\omega^2\epsilon_1$, whose coordinate sum is nonzero. It also does not preserve the edge-sum line. Thus this order column is not already a tame descent of the earlier cellular chain complex or its $H_1$.

For the split local coefficient $\pi=a+b\omega$, the natural rotation polynomial on the balanced carrier is $F=aI+bC$. Since $g$ fixes $a,b\in k$, its commutator is
\[
 \Theta F-F\Theta=b(\omega^2-1)C\Theta.
\]
It is nonzero when $b\ne0$. In particular both prime-seven twists $(a,b)=(3,2),(-1,-3)$ fail this test. Prime powers with $b=0$ are not excluded by this particular calculation; that exceptional commutation does not prove a global comparison.

Replacing $C$ by $vC$ gives an invariant carrier map, whose matrix on $L$ is $aI+bA$. Direct determinant expansion gives
\[
 \det(1-T(aI+bA))=(1-aT)^3-ub^3T^3.
\]
At $u=0$ this is $(1-aT)^3$, not $P_d(T)$, which has degree two. The weighted replacement therefore does not recover the established local factor. Neither $C$ nor $vC$ is generally $\Lambda$-linear: they intertwine its rotation. Their sums with the identity have not thereby become endomorphisms in the ordinary order-module category. They are legitimate carrier tests, with exactly the failures just exhibited.

\section{The full orbit is the normalization, and arbitrary kernels still fail}
Twist the tame action by the three characters: $\Theta_r=\omega^r\Theta$, $r\in\mathbb Z/3$. Its fixed lattice is
\[
 L_r=\bigoplus_i Rv^{[i-r]_3}\epsilon_i,
\]
with the cut at arrow $r$. Rotation gives an actual $R$-linear isomorphism $C:L_r\to L_{r+1}$, taking each displayed basis vector to the corresponding next one. Thus the direct sum of all three cuts has a strict rotation, at the cost of vertex rank three. Before invariants the orbit is $S$-rank nine and underlying $R$-rank twenty-seven; its invariant lattice is $R$-rank nine. The original orbit $S$-lattice still has nontrivial inertia at zero and need not be the base change of these invariants.

There is an explicit integral identification of the invariant orbit. In the standard hereditary order $B_+$, the $j$-th column has entries $u^{t_{ij}}R$ with $t_{ij}=1$ for $i<j$ and zero otherwise. In its natural column basis the arrow coefficient has exponent
\[
 \delta_{i0}+t_{i-1,j}-t_{ij}=\delta_{ij}.
\]
Its cut is therefore exactly $j$, identifying $B_+e_j$ with $L_j$. Consequently $\bigoplus_rL_r\simeq B_+$ as left $\Lambda_R$-modules. Adding the exchanged minus-branch orbit gives the hereditary normalization $\widetilde\Lambda=B_+\times B_-$, with full deck action and vertex rank six. On a basic projective with external arithmetic matrix this yields $P_d^6$.

In fact this normalization replacement restores the \emph{ordinary determinant for every linear endomorphism}. For a finite projective $P$, the inclusion $P\to P\otimes_\Lambda\widetilde\Lambda$ becomes an isomorphism over $K=k((u))$, since $\widetilde\Lambda_K=\Lambda_K$. Its extended endomorphism and the original $f$ therefore have the same generic characteristic determinant. Both polynomials have coefficients in $R$, which injects into $K$, so they agree already over $R$. The bounded-complex ratios also agree. This symmetry repair returns the previously refuted ordinary determinant, not merely its external tensor exponent.

For completeness we now prove Theorem~\ref{thm:full-deck-kernel} for arbitrary finite free kernels, without assuming they split integrally into branches or are projective over the order. Let
\[
 p_\pm=u^{-1}\sum_i c_i^\pm\in\Lambda_K.
\]
These are central generic branch idempotents. Put $W_\pm=W_0\cap p_\pm(W_0\otimes_R K)$. They are finite free $R$-lattices spanning those generic summands, stable under $\Lambda_R$. Rotation fixes them, branch swap exchanges them. Their vertex ranks are constant on each branch: the arrow maps become invertible over $K$ because a full cycle acts by $u$. Write these ranks $r_+,r_-$. No equality $W_0=W_+\oplus W_-$ over $R$ has been assumed.

On the plus branch choose integral vertex bases. The three square arrow matrices $X_0,X_1,X_2$ have product $uI_{r_+}$ around every cycle. Thus
\[
 \sum_i\operatorname{val}_u(\det X_i)=r_+.
\]
The rotation lift is an invertible matrix between each pair of vertex lattices, hence its determinant is an $R$-unit. Its intertwining identities make all three arrow determinant valuations equal. Therefore $r_+=3m_+$; similarly $r_-=3m_-$. The branch-swap lift gives $r_+=r_-$, so generic rank additivity gives $\operatorname{rank}_R e_iW_0=6m$. Only invertible intertwiners were used; imposing the full deck group relations is stronger than necessary.

For the general determinant statement reduce modulo $u$ and filter each vertex by $e_iJ^a(W_0/uW_0)$. Since $J^4\subset u\Lambda$, this is a finite filtration. In $\Phi_{W_0}(P)=\bigoplus_i(e_iW_0)^{n_i}$, positive-length entries of $f$ raise the filtration, while the degree-zero matrix $A_i$ repeats on its $i$-th graded pieces. Their total dimension is $6m$. Multiplying the graded characteristic determinants proves the asserted power for \emph{every} $f$. Flatness of the kernel as a left $\Lambda$-module was not used; bounded right projectives compute the derived functor and give finite free $R$-outputs. Denominator constant terms one permit specialization of every bounded alternating ratio. The regular left kernel $\Lambda_R$ attains $m=1$, making the rank bound sharp.

The remaining arithmetic problem therefore requires a different hypothesis: a genuinely twisted determinant target, a specified non-linear/semilinear correspondence with its own comparison, or an infinite analytic realization. One cannot continue searching inside this closed finite-free, fixed-center, full-deck kernel class for a missing small estimate. The local purity and descent computations supply no global positive spectral limit. Moreover $k/\mathbb Q$ has degree two: ordinary restriction of scalars sends a rational local factor $P_d$ to $P_d^2$. A relative $k$-determinant must not silently be identified with the former rational determinant. Finite-field curve zeta, global elliptic $L$-functions and Dirichlet GRH keep their distinct scopes and centers.



\chapter{Literal formal and computational coverage}
\texttt{DessinNodeOrder.lean} defines the fiber-product subring of two genuine power-series rings. Its twenty printed declarations establish diagonal injectivity, nonzero branches, branch zero products, factorizations and annihilator identities, regular $u$, augmentation surjectivity and exactness, bijective power-series coordinates, square-zero alternating maps, three kernel-equals-range statements and projectivity of the actual regular/product terms. The coordinate statement is a bijection; the file does not yet package it as an $R$-linear equivalence or declare the general order freeness theorem.

\texttt{DessinInvariantComponents.lean} has ten printed declarations: three CK1-collapse components, determinant change of bases with an actual inverse, symmetric-square zero, identity, composition and relation preservation, the degree-two Schur character identity, and the nodal central-parameter obstruction. Its symmetric-square matrix is
\[
 \operatorname{Sym}^2\!\begin{pmatrix}a&b\\c&d\end{pmatrix}
 =\begin{pmatrix}a^2&ab&b^2\\2ac&ad+bc&2bd\\c^2&cd&d^2\end{pmatrix}.
\]
This is the ordinary monomial basis convention; composition is proved universally over any commutative ring. The full project suite prints types and foundational dependencies; no theorem about general LR decompositions, abstract derived equivalences, C* $K$-theory or GRH is inferred from these component proofs.

The exact Python checks use fixed rational matrices and exhaustive finite sets, with no random seeds. They construct the displayed permutations and paths, evaluate rational Cartan identities, check formal Euler products through a stated degree, verify unrestricted-output Fock basis identities on bounded input sets, compute incidence Smith data, construct all nine string representations and signed maps, and test Young/LR/Cauchy examples and the seven-element-field counts. Those diagnostics are reproducible finite corroboration. The universal infinite claims rest on the proofs or explicitly cited theorems, and the arithmetic comparison remains open.

The ten additional \texttt{DessinPullbackComponents.lean} statements use literal two-state matrix units and a subring of a product of arbitrary rings, with equal images under two ring homomorphisms. The section hypothesis of its projection-surjectivity theorem is explicit. It does not construct a topological C*-pullback or assert that a commutative diagram is automatically a pullback. The no-residue theorem is universal algebra, with a nontrivial commutative target; it is not an assumed arithmetic positivity premise.

\texttt{ThetaJacobi.lean} adds sixteen declarations, with actual finite fields and MulChar rather than a hypothetical arithmetic kernel. Its norm theorem invokes mathlib's proved Jacobi inverse-product identity and derives characteristic inequality from finiteness; all three character nontrivialities are explicit. Other declarations verify the rescaling substitution, CM multiplication/polynomial/Gram positivity components, exact prime-seven matrix polynomials, and actual 2-by-3 boundary/3-by-2 cellular inclusion, homology restriction, three-edge determinant and rational vertex chain-map equations. It does not formalize elliptic point counting or an order-to-Frobenius functor.

The final three \texttt{ThetaJacobi.lean} declarations verify the actual conjugation matrix identities and the real two-dimensional inert-prime determinant obstruction. They do not assume any unknown zero location or exclude complex/enlarged spectral realizations.


\texttt{ThetaProjectiveMultiplicity.lean} adds five declarations. Arbitrary upper-block matrices with six identical diagonal blocks have the sixth-power determinant and characteristic-factor formulas. A literal power-series constant-coefficient map specializes matrix determinants, for each scalar parameter; the coefficientwise polynomial/order interface remains prose. A nonzero quadratic cannot be a rational-function sixth power, proved by polynomial degrees. Three genuine power series with product $X$ cannot have the cyclically equivalent residue-vanishing conditions forced by invertible rank-one arrow intertwiners. The derivation of the path blocks from the full theta order, projective decomposition, derived branch functor and deck intertwiners is not fully formalized. The complete analytic suite now has208printed declarations; no hidden arithmetic positivity assumption occurs.


\texttt{ThetaTameDescent.lean} adds twelve printed component declarations (and one auxiliary polynomial lemma). They use genuine formal-power-series rescaling and its coefficient operator, proving the coefficient action, invariant-coefficient vanishing and the tame operator's cube. Actual three-state matrices satisfy the clock/rotation commutator, comparison determinant and closed-fiber equations, weighted-cycle cube and full characteristic polynomial, noncommutation of a nonconstant rotation polynomial, trivial whole-column determinant actions, and the same-operator comparison for $vC$. The classification of all fixed series by residues, full order action and hereditary invariant algebra, stack interpretation, valuation/rank theorem, general coefficient-functor filtration and GRH are not thereby fully formalized. No unknown positivity premise is assumed.
